\documentclass[11pt]{article}
\usepackage[margin=1in]{geometry}
\usepackage[T1]{fontenc}
\usepackage[utf8]{inputenc}
\usepackage[french]{babel}
\usepackage{amsmath,amssymb,amsthm}
\usepackage[colorlinks=true,linkcolor=blue,urlcolor=blue]{hyperref}
\newtheorem{problem}{Problème}
\newenvironment{refs}{\par\small\noindent\emph{Références :}\begin{itemize}\setlength\itemsep{0pt}}{\end{itemize}\normalsize}
\title{Hors de la Caverne \\ \Large Probabilités}
\author{}
\date{}
\begin{document}
\maketitle
\noindent\emph{Des problèmes, pas de réponses.} Chaque problème ci-dessous est posé
sans solution. Les références indiquent où trouver les connaissances nécessaires.
\medskip


\section*{Collège}

\begin{problem}[{Sept contre douze --- Collège}]
\textbf{PRB-45.} On lance deux dés équilibrés à six faces, l'un rouge et l'un vert, et on
additionne les deux résultats. Un joueur affirme : \og{} il y a onze sommes possibles, de 2 à 12, donc
chacune a une chance sur onze de sortir \fg{}.
\begin{enumerate}
\item[(a)] Expliquez pourquoi ce sont les $ 36 $ couples (dé rouge, dé vert) qui forment les cas
également possibles, et non les onze sommes. Votre argument doit dire précisément ce qui distingue
les deux façons de compter.
\item[(b)] Dressez le tableau des 36 issues, puis comparez la probabilité d'obtenir 7 et celle d'obtenir
12. Justifiez l'écart que vous trouvez en décrivant, en français et sans formule, ce que la somme 7
a de particulier parmi toutes les sommes.
\item[(c)] Le joueur se défend : \og{} j'ai lancé les dés 36 fois et je n'ai obtenu 7 que quatre fois ; votre
calcul est donc faux \fg{}. Rédigez une réponse qui dit ce qu'une probabilité promet et ce qu'elle ne
promet pas.
\end{enumerate}

\end{problem}

\noindent Le grand-duc de Toscane posa à Galilée, vers 1620, la même question pour trois
dés : les joueurs savaient d'expérience que 10 sortait plus souvent que 9, alors que les deux
sommes s'écrivent chacune de six façons. Galilée trancha en quelques lignes, en comptant les issues
ordonnées plutôt que les descriptions d'issues --- c'est le premier raisonnement de probabilité qui
mérite ce nom, et c'est exactement le geste demandé à la question (a). La version à trois dés est
PRB-01, sur cette même page.

\begin{refs}
\item Contexte : Probabilité --- Wikipédia (\url{https://fr.wikipedia.org/wiki/Probabilit%C3%A9})
\item Contexte : Dé --- Wikipédia (\url{https://fr.wikipedia.org/wiki/D%C3%A9})
\item Contexte : Galilée --- Wikipédia (\url{https://fr.wikipedia.org/wiki/Galil%C3%A9e_(savant)})
\end{refs}

\bigskip

\begin{problem}[{La suite qui n'a pas l'air au hasard --- Collège}]
\textbf{PRB-46.} Deux élèves devaient produire une suite de 30 résultats de pile ou face. L'un a
réellement lancé une pièce et noté ce qu'il obtenait ; l'autre a inventé sa suite de tête. Voici
les deux copies, où P désigne pile et F face.

\noindent Suite 1 : P F P F F P F P P F P F F P F P F P P F P F P F F P F P F P 
Suite 2 : P P F F F F P F P P P F F P F F F F F P P F P P P P F P F F
\begin{enumerate}
\item[(a)] Pour chacune des deux suites, calculez la fréquence de \og{} pile \fg{} et relevez la plus longue série
de résultats identiques consécutifs. Présentez vos résultats dans un tableau.
\item[(b)] Dites laquelle des deux suites vous croyez inventée et défendez votre choix. On attend un
argument appuyé sur les nombres de la question (a), pas une impression visuelle.
\item[(c)] Justifiez qu'une suite précise de 30 lancers a exactement la même probabilité qu'une autre
suite précise de 30 lancers --- y compris celle qui donne trente fois \og{} pile \fg{}. Expliquez ensuite,
en un paragraphe rédigé, pourquoi cela ne contredit pas votre réponse à la question (b).
\end{enumerate}

\end{problem}

\noindent Cette expérience est une démonstration de cours devenue classique --- souvent
attribuée au mathématicien américain Ted Hill : le professeur sort de la salle, une moitié de la
classe lance vraiment une pièce cent fois, l'autre moitié invente sa suite, et il retrouve les
tricheurs en quelques secondes. Les êtres humains alternent trop et produisent presque toujours
des séries trop courtes, parce que nous confondons \og{} au hasard \fg{} et \og{} bien mélangé \fg{}. La question
(c) est le c\oe{}ur de l'affaire : le hasard ne rend pas certaines suites impossibles, il rend rares
certaines \textit{propriétés} des suites --- la version chiffrée de cet argument est PRB-05, dans la
bande Lycée.

\begin{refs}
\item Contexte : Pile ou face --- Wikipédia (\url{https://fr.wikipedia.org/wiki/Pile_ou_face})
\item Contexte : Hasard --- Wikipédia (\url{https://fr.wikipedia.org/wiki/Hasard})
\item Contexte : Ted Hill --- Wikipédia (en anglais) (\url{https://en.wikipedia.org/wiki/Ted_Hill_(mathematician)})
\end{refs}

\bigskip

\begin{problem}[{Deux chaussettes dans le noir --- Collège, short}]
\textbf{PRB-47.} Un tiroir contient six chaussettes noires et quatre
chaussettes bleues mélangées, et vous en tirez deux au hasard dans l'obscurité, sans remettre la
première avant de tirer la seconde. Déterminez la probabilité d'obtenir deux chaussettes de la même
couleur en justifiant chaque branche de votre arbre, puis expliquez précisément ce qui changerait
si vous remettiez la première chaussette dans le tiroir.
\end{problem}

\noindent Le tiroir de chaussettes est un exercice de manuel depuis au moins les années
1950, et il circule sous deux formes que tout oppose. La forme \og{} combien faut-il en tirer pour être
\textit{certain} d'avoir une paire ? \fg{} ne demande aucune probabilité : elle relève du principe des
tiroirs et se règle par un décompte. La forme posée ici demande un arbre et donne un nombre entre 0
et 1. Savoir laquelle des deux questions on est en train de résoudre est déjà une compétence de
mathématicien.

\begin{refs}
\item Contexte : Problème d'urne --- Wikipédia (\url{https://fr.wikipedia.org/wiki/Probl%C3%A8me_d%27urne})
\item Contexte : Arbre de probabilité --- Wikipédia (\url{https://fr.wikipedia.org/wiki/Arbre_de_probabilit%C3%A9})
\item Contexte : Principe des tiroirs --- Wikipédia (\url{https://fr.wikipedia.org/wiki/Principe_des_tiroirs})
\end{refs}

\bigskip

\begin{problem}[{La roue du forain --- Collège}]
\textbf{PRB-48.} À la fête foraine, une roue est partagée en huit secteurs de même taille,
numérotés de 1 à 8 ; le forain jure qu'elle est parfaitement équilibrée. Vous restez une heure à
noter les résultats. Sur 400 tours, les secteurs 1 à 8 sont sortis respectivement 62, 54, 47, 51,
58, 49, 66 et 13 fois.
\begin{enumerate}
\item[(a)] Calculez la fréquence de sortie de chaque secteur et l'effectif que l'on attendrait pour
chacun si la roue était équilibrée. Rassemblez le tout dans un tableau et commentez ligne par
ligne.
\item[(b)] Le forain répond : \og{} le hasard fait ce qu'il veut, vos écarts ne prouvent rien \fg{}. Il a raison
sur le principe. Construisez malgré tout l'argument le plus solide que vous puissiez écrire au
collège pour soutenir que la roue est truquée --- puis dites honnêtement ce qui manque à votre
argument pour être une preuve.
\item[(c)] Le secteur 8 est celui qui fait gagner le gros lot. Expliquez pourquoi un forain malhonnête a
intérêt à truquer précisément ce secteur-là, et pourquoi un joueur qui ne regarde qu'une dizaine de
tours n'a pratiquement aucune chance de s'en apercevoir.
\end{enumerate}

\end{problem}

\noindent Cette méthode a réellement fait fortune. En 1873, l'ingénieur anglais Joseph
Jagger paya des employés pour relever pendant des semaines les numéros sortis aux roulettes du
casino de Monte-Carlo ; l'une des roues était légèrement déséquilibrée, certains numéros sortaient
trop souvent, et Jagger repartit avec une somme considérable. Les casinos rééquilibrent depuis
leurs roues et changent leur emplacement, précisément pour empêcher ce genre de relevé. L'idée que
la fréquence observée finit par ressembler à la probabilité porte un nom --- la loi des grands
nombres --- et se démontre au lycée ; ce qui vous manque ici, et que la question (b) vous demande de
nommer, est justement l'outil qui dit à partir de quel écart on peut cesser de croire au
forain.

\begin{refs}
\item Contexte : Fréquence (statistiques) --- Wikipédia (\url{https://fr.wikipedia.org/wiki/Fr%C3%A9quence_(statistiques)})
\item Contexte : Loi des grands nombres --- Wikipédia (\url{https://fr.wikipedia.org/wiki/Loi_des_grands_nombres})
\item Contexte : Joseph Jagger --- Wikipédia (en anglais) (\url{https://en.wikipedia.org/wiki/Joseph_Jagger})
\end{refs}

\bigskip

\begin{problem}[{Le jeu est-il équitable ? --- Collège}]
\textbf{PRB-49.} Alice et Bruno lancent ensemble deux dés équilibrés à six faces. Dans chacune
des règles ci-dessous, il n'y a ni match nul ni relance.
\begin{enumerate}
\item[(a)] Règle 1 : Alice gagne si la somme des deux dés est paire, Bruno si elle est impaire. Le jeu
est-il équitable ? Justifiez à partir du tableau des 36 issues.
\item[(b)] Règle 2 : Alice gagne si le produit des deux dés est pair, Bruno s'il est impair. Justifiez
que ce jeu n'est pas équitable, et dites de combien exactement il penche.
\item[(c)] Pour compenser la règle 2, Bruno propose de recevoir $ k $ euros chaque fois qu'il gagne et
de payer 1 euro chaque fois qu'il perd. Déterminez, en raisonnant sur un grand nombre de parties,
la valeur de $ k $ qui rend ce jeu équitable, et justifiez-la. Expliquez ensuite pourquoi
\og{} équitable \fg{} ne veut pas dire la même chose à la question (a) et à la question (c).
\end{enumerate}

\end{problem}

\noindent La notion de jeu équitable est le berceau de toute la théorie : en 1654, le
chevalier de Méré --- de son vrai nom Antoine Gombaud --- apporta à Blaise Pascal des questions de ce
type, et la correspondance entre Pascal et Fermat qui s'ensuivit fonda la discipline. Huygens en
tira en 1657 le premier traité imprimé, où il définit la valeur d'un jeu comme la mise qu'un joueur
devrait accepter de payer pour y entrer : c'est l'ancêtre direct de votre question (c). L'outil qui
la formalise s'appelle l'espérance et se rencontre au lycée ; au collège, le raisonnement \og{} sur un
grand nombre de parties \fg{} suffit, et il a l'avantage de dire ce que l'espérance signifie avant de
dire comment on la calcule.

\begin{refs}
\item Contexte : Antoine Gombaud, chevalier de Méré --- Wikipédia (\url{https://fr.wikipedia.org/wiki/Antoine_Gombaud})
\item Contexte : Christian Huygens --- Wikipédia (\url{https://fr.wikipedia.org/wiki/Christian_Huygens})
\item Contexte : Espérance mathématique --- Wikipédia (\url{https://fr.wikipedia.org/wiki/Esp%C3%A9rance_math%C3%A9matique})
\end{refs}

\bigskip

\begin{problem}[{Un anniversaire commun dans votre classe --- Collège}]
\textbf{PRB-50.} Un élève d'une classe de 30 affirme : \og{} nous sommes 30 et l'année compte 365
jours, donc il n'y a presque aucune chance que deux d'entre nous soient nés le même jour \fg{}.
\begin{enumerate}
\item[(a)] Comptez le nombre de paires d'élèves que l'on peut former dans cette classe. Expliquez
pourquoi c'est ce nombre-là, et non 30, qu'il faut comparer à 365, et pourquoi le raisonnement de
l'élève passe à côté du problème.
\item[(b)] Distinguez soigneusement deux questions : \og{} un autre élève est-il né le même jour que
\textit{moi} ? \fg{} et \og{} deux élèves quelconques de la classe sont-ils nés le même jour ? \fg{}. Justifiez
que la seconde a beaucoup plus de chances d'être vraie que la première, sans calculer aucune des
deux probabilités.
\item[(c)] Menez l'expérience. Relevez les dates de naissance de votre classe et de trois classes
voisines, comptez combien de ces quatre classes contiennent au moins un anniversaire commun, et
confrontez le résultat à votre raisonnement de la question (a). Discutez enfin ce que quatre
classes permettent --- et ne permettent pas --- de conclure.
\end{enumerate}

\end{problem}

\noindent Le paradoxe des anniversaires fut posé sous cette forme par Richard von Mises en
1939 et rendu célèbre par le manuel de William Feller. Il n'a rien d'un paradoxe logique : c'est
une contradiction entre un calcul juste et une intuition fausse, et l'intuition se trompe parce
qu'elle compte les personnes là où il faut compter les paires. Le calcul exact, avec le seuil
fameux de 23 personnes, est PRB-03 dans la bande Lycée ; la question (a) ci-dessus en contient déjà
toute l'idée, et la question (c) vous fait faire ce que von Mises n'a jamais pris la peine de
faire.

\begin{refs}
\item Contexte : Paradoxe des anniversaires --- Wikipédia (\url{https://fr.wikipedia.org/wiki/Paradoxe_des_anniversaires})
\item Contexte : Probabilité --- Wikipédia (\url{https://fr.wikipedia.org/wiki/Probabilit%C3%A9})
\end{refs}

\bigskip

\begin{problem}[{Deux enfants, dont une fille --- Collège}]
\textbf{PRB-51.} On admet qu'à chaque naissance, garçon et fille sont également probables, et
on considère une famille de deux enfants.
\begin{enumerate}
\item[(a)] Dressez l'arbre des quatre familles possibles, en distinguant l'aîné du cadet, et justifiez
que ces quatre issues sont également probables.
\item[(b)] On vous dit seulement : \og{} cette famille a au moins une fille \fg{}. Déterminez la probabilité que
les deux enfants soient des filles, en justifiant à partir de votre arbre quelles branches vous
gardez et lesquelles vous éliminez.
\item[(c)] Situation différente : vous croisez dans la rue un enfant de cette famille, choisi au hasard
parmi les deux, et c'est une fille. Reprenez l'arbre --- il faut maintenant y faire figurer le choix
de l'enfant que vous rencontrez --- et déterminez à nouveau la probabilité que les deux soient des
filles.
\item[(d)] Les deux réponses ne sont pas les mêmes. Expliquez, dans un paragraphe rédigé, ce qui a changé
entre (b) et (c), et pourquoi une énigme de probabilité n'est pas encore un problème tant qu'on n'a
pas dit \textit{comment} l'information a été obtenue.
\end{enumerate}

\end{problem}

\noindent Martin Gardner posa l'énigme dans sa chronique de \textit{Scientific American} en
1959, publia une réponse, puis fit machine arrière : la question, telle qu'il l'avait écrite, ne
déterminait pas la réponse. Ce n'est pas une faiblesse de l'énoncé, c'est le contenu même du
problème --- en probabilité, ce qui compte n'est pas le fait que vous apprenez mais la procédure qui
vous l'a fait apprendre. La version calculée avec des probabilités conditionnelles est PRB-23,
dans la bande Lycée ; ici, l'arbre suffit, et il rend la morale plus visible encore.

\begin{refs}
\item Contexte : Paradoxe des deux enfants --- Wikipédia (\url{https://fr.wikipedia.org/wiki/Paradoxe_des_deux_enfants})
\item Contexte : Arbre de probabilité --- Wikipédia (\url{https://fr.wikipedia.org/wiki/Arbre_de_probabilit%C3%A9})
\item Contexte : Martin Gardner --- Wikipédia (\url{https://fr.wikipedia.org/wiki/Martin_Gardner})
\end{refs}

\bigskip

\begin{problem}[{Trois dés qui tournent en rond --- Collège}]
\textbf{PRB-52.} Trois dés équilibrés à six faces portent les valeurs
$ A = (2,2,4,4,9,9) $, $ B = (1,1,6,6,8,8) $ et $ C = (3,3,5,5,7,7) $. Deux joueurs choisissent
chacun un dé, les lancent, et celui qui obtient le plus grand nombre gagne la partie.
\begin{enumerate}
\item[(a)] Pour chacun des trois duels --- A contre B, B contre C, puis C contre A --- dressez le tableau des
36 issues et déterminez, en justifiant, quel dé a le plus de chances de l'emporter.
\item[(b)] Comparez vos trois réponses. Énoncez avec précision l'idée reçue qu'elles mettent en défaut,
et expliquez en français pourquoi cette idée reçue semblait pourtant évidente.
\item[(c)] Vous proposez ce jeu à un camarade en le laissant choisir son dé le premier. Expliquez
pourquoi vous avez alors l'avantage, et dites ce qui change si c'est vous qui devez choisir en
premier.
\end{enumerate}

\end{problem}

\noindent Ces dés sont dus au statisticien Bradley Efron et furent rendus célèbres par
Martin Gardner dans \textit{Scientific American} en 1970. Leur intérêt n'est pas l'arithmétique mais
ce qu'ils détruisent : presque tout le monde suppose en silence que \og{} a plus de chances de gagner
contre \fg{} se transmet de proche en proche, comme \og{} est plus grand que \fg{} se transmet chez les
nombres. Cette hypothèse tacite sous-tend les classements sportifs, les comparatifs de produits et
certains modes de scrutin --- et voici un contre-exemple qui tient dans trois tableaux. La version à
quatre dés est PRB-24, dans la bande Lycée.

\begin{refs}
\item Contexte : Dés non transitifs --- Wikipédia (\url{https://fr.wikipedia.org/wiki/D%C3%A9s_non_transitifs})
\item Contexte : Martin Gardner --- Wikipédia (\url{https://fr.wikipedia.org/wiki/Martin_Gardner})
\end{refs}

\bigskip

\begin{problem}[{\og{} Deux issues, donc une chance sur deux \fg{} --- Collège, short}]
\textbf{PRB-53.} Un élève affirme : \og{} quand je joue au loto, soit je gagne,
soit je perds ; il n'y a que deux issues, donc j'ai une chance sur deux de gagner \fg{}. Débusquez
l'erreur en énonçant précisément l'hypothèse qui manque, puis construisez deux expériences n'ayant
chacune que deux issues, l'une où \og{} une chance sur deux \fg{} est juste et l'autre où c'est faux, et
justifiez les deux cas.
\end{problem}

\noindent L'erreur a une lettre de noblesse. Dans l'article \og{} Croix ou pile \fg{} de
l'\textit{Encyclopédie} (1754), d'Alembert soutient qu'en deux lancers d'une pièce la probabilité
d'obtenir au moins une fois croix vaut $ 2/3 $, parce qu'il ne voit que trois cas : croix au
premier coup, pile puis croix, pile puis pile. Il compte des descriptions d'issues au lieu de
compter des issues également probables --- la faute même de l'élève ci-dessus, commise par l'un des
plus grands mathématiciens du siècle. L'hypothèse d'équiprobabilité n'est jamais gratuite : elle se
justifie ou elle se démontre, jamais elle ne se suppose au vu du nombre de cases.

\begin{refs}
\item Contexte : Équiprobabilité --- Wikipédia (\url{https://fr.wikipedia.org/wiki/%C3%89quiprobabilit%C3%A9})
\item Contexte : Jean Le Rond d'Alembert --- Wikipédia (\url{https://fr.wikipedia.org/wiki/Jean_Le_Rond_d%27Alembert})
\item Contexte : Loterie --- Wikipédia (\url{https://fr.wikipedia.org/wiki/Loterie})
\end{refs}

\bigskip

\begin{problem}[{Deux urnes et un dé --- Collège}]
\textbf{PRB-54.} L'urne n$^\circ$ 1 contient 3 boules rouges et 2 boules vertes ; l'urne n$^\circ$ 2 contient
1 boule rouge et 4 boules vertes. On lance un dé équilibré : si l'on obtient 5 ou 6, on tire une
boule dans l'urne n$^\circ$ 1, sinon on tire une boule dans l'urne n$^\circ$ 2.
\begin{enumerate}
\item[(a)] Dressez l'arbre de cette expérience en portant sur chaque branche la probabilité
correspondante, et justifiez que la somme des probabilités des branches issues d'un même n\oe{}ud vaut
toujours 1.
\item[(b)] Déterminez la probabilité de tirer une boule rouge, en justifiant chaque étape de votre
calcul sur l'arbre.
\item[(c)] Un camarade objecte : \og{} il y a en tout 4 boules rouges parmi les 10 boules des deux urnes,
donc la probabilité de tirer une rouge vaut $ 4/10 $ \fg{}. Expliquez pourquoi ce raisonnement est
faux ici, puis décrivez précisément une règle de tirage --- différente de celle de l'énoncé --- pour
laquelle il serait juste.
\end{enumerate}

\end{problem}

\noindent Les urnes remplies de boules de couleur sont le laboratoire de la théorie des
probabilités depuis Jacques Bernoulli, dont l'\textit{Ars conjectandi} (publié en 1713) en fait
l'instrument de base du raisonnement sur le hasard. Leur intérêt est d'être totalement
transparentes : on voit ce qu'il y a dedans, il ne reste donc rien à deviner et tout à démontrer.
La question (c) contient la seule difficulté véritable --- une probabilité n'est une fraction de cas
favorables que lorsque les cas sont également probables, et deux urnes tirées avec des chances
inégales ne se laissent pas verser dans un même sac.

\begin{refs}
\item Contexte : Problème d'urne --- Wikipédia (\url{https://fr.wikipedia.org/wiki/Probl%C3%A8me_d%27urne})
\item Contexte : Arbre de probabilité --- Wikipédia (\url{https://fr.wikipedia.org/wiki/Arbre_de_probabilit%C3%A9})
\end{refs}

\bigskip

\begin{problem}[{Au moins un six --- Collège, short}]
\textbf{PRB-55.} Un élève raisonne ainsi sur le lancer de deux dés
équilibrés : \og{} la probabilité d'obtenir un six sur un dé vaut $ 1/6 $, donc avec deux dés elle
vaut $ 1/6 + 1/6 = 1/3 $ \fg{}. Déterminez la probabilité d'obtenir au moins un six en la justifiant
sur le tableau des 36 issues, puis énoncez la condition sous laquelle additionner deux probabilités
est légitime et montrez qu'elle n'est pas remplie ici.
\end{problem}

\noindent L'addition naïve des probabilités est l'erreur la plus répandue de tout
l'enseignement de la discipline, et elle a résisté longtemps : les premiers auteurs --- Cardan au
XVIe siècle dans son \textit{Liber de ludo aleae}, Galilée vers 1620 --- ont dû forger la
notion de \og{} cas également possibles \fg{} précisément pour se donner un garde-fou contre ce genre de
glissement. Le point à retenir dépasse les dés : additionner des probabilités revient à supposer
que deux choses ne peuvent pas arriver ensemble, et c'est une hypothèse sur le monde, pas une règle
de calcul.

\begin{refs}
\item Contexte : Événement (probabilités) --- Wikipédia (\url{https://fr.wikipedia.org/wiki/%C3%89v%C3%A9nement_(probabilit%C3%A9s)})
\item Contexte : Jérôme Cardan --- Wikipédia (\url{https://fr.wikipedia.org/wiki/J%C3%A9r%C3%B4me_Cardan})
\end{refs}

\bigskip

\begin{problem}[{La pièce qui \og{} doit \fg{} tomber sur pile --- Collège, short}]
\textbf{PRB-56.} Une pièce équilibrée vient de tomber cinq fois de suite sur
face, et un joueur en conclut : \og{} pile est en retard, il a maintenant plus d'une chance sur deux de
sortir au prochain lancer \fg{}. Réfutez ce raisonnement, puis expliquez comment il peut être faux
alors même qu'il est vrai que, sur un très grand nombre de lancers, la fréquence de \og{} pile \fg{} se
rapproche de $ 1/2 $.
\end{problem}

\noindent Laplace décrit déjà cette illusion dans son \textit{Essai philosophique sur les
probabilités} (1814) : il y raconte des joueurs de loterie qui misent sur les numéros non sortis
depuis longtemps, persuadés qu'ils sont \og{} dus \fg{}, et il en fait l'exemple type des erreurs de
jugement sur le hasard. Deux siècles plus tard, les sites de loterie publient encore la liste des
numéros \og{} en retard \fg{}, et les joueurs les jouent. La seconde partie de la question est la vraie :
la fréquence se rapproche bien de $ 1/2 $, mais elle y arrive en \textit{diluant} les écarts
passés, jamais en les corrigeant --- la version chiffrée de cette distinction est PRB-05, dans la
bande Lycée.

\begin{refs}
\item Contexte : Erreur du parieur --- Wikipédia (\url{https://fr.wikipedia.org/wiki/Erreur_du_parieur})
\item Contexte : Pierre-Simon de Laplace --- Wikipédia (\url{https://fr.wikipedia.org/wiki/Pierre-Simon_de_Laplace})
\item Contexte : Loi des grands nombres --- Wikipédia (\url{https://fr.wikipedia.org/wiki/Loi_des_grands_nombres})
\end{refs}

\bigskip


\section*{Lycée}

\begin{problem}[{Les dés de Galilée --- Lycée}]
\textbf{PRB-01.} On lance trois dés équilibrés. Les sommes 9 et 10 s'écrivent
chacune comme somme de trois valeurs de dés d'exactement six façons
($ 9 = 1{+}2{+}6 = 1{+}3{+}5 = \dots $, et de même pour 10) --- ce qui laisse croire qu'elles
devraient être équiprobables. Les joueurs savaient d'expérience que le 10 sort plus souvent.
Donnez raison aux joueurs : construisez l'univers correct des $ 6^3 = 216 $ issues ordonnées
équiprobables, montrez que
$ P(\text{somme}=10) = \frac{27}{216} > \frac{25}{216} = P(\text{somme}=9) $, et expliquez
exactement où l'argument des \og{} six façons pour chacune \fg{} se trompe. Déduisez-en ensuite la loi
complète de la somme de \textit{deux} dés et démontrez que 7 en est la valeur la plus probable.
\end{problem}

\noindent Le grand-duc de Toscane soumit l'énigme du 9 contre le 10 à Galilée, qui la
trancha dans une courte note (\textit{Sopra le scoperte dei dadi}, vers 1620) ; Cardan était
parvenu à des idées voisines dans son \textit{Liber de ludo aleae} un demi-siècle plus tôt, resté
inédit jusqu'en 1663. La leçon --- compter les issues \textit{équiprobables}, et non les descriptions
d'issues --- est la première idée sérieuse de la théorie des probabilités, et les joueurs de dés
avaient relevé empiriquement un écart de 1/108 avant que les mathématiques ne sachent
l'expliquer.

\begin{refs}
\item Contexte : Histoire des probabilités --- Wikipédia (\url{https://fr.wikipedia.org/wiki/Histoire_des_probabilit%C3%A9s})
\item Lecture : William Feller, \textit{An Introduction to Probability Theory and Its Applications}, vol. 1, ch. I
\item Cours : MIT OCW 18.600 Probability and Random Variables (\url{https://ocw.mit.edu/courses/18-600-probability-and-random-variables-fall-2019/})
\end{refs}

\bigskip

\begin{problem}[{Les deux paris du Chevalier --- Lycée}]
\textbf{PRB-02.} Le chevalier de Méré gagnait de l'argent en pariant qu'il
obtiendrait au moins un six en quatre lancers d'un dé, puis en perdit en pariant qu'il obtiendrait
au moins un double-six en vingt-quatre lancers de deux dés --- alors qu'il tenait les deux paris
pour équivalents (\og{} 4 est à 6 ce que 24 est à 36 \fg{}).
\begin{enumerate}
\item[(i)] Démontrez que le premier pari est favorable :
$ 1 - (5/6)^4 > 1/2 $.
\item[(ii)] Démontrez que le second ne l'est pas :
$ 1 - (35/36)^{24} < 1/2 $, et déterminez, avec démonstration, le plus petit nombre de lancers
qui rend favorable le pari sur le double-six.
\item[(iii)] Mettez le doigt sur l'étape fausse du
raisonnement de proportionnalité.
\item[(iv)] Le problème des partis : deux joueurs misent la même somme sur un jeu de pile ou face
équilibré où le premier à remporter un nombre fixé de manches emporte tout ; ils doivent
s'arrêter alors qu'il manque encore 2 victoires au joueur A et 3 au joueur B. Déterminez, avec
démonstration, le partage équitable des mises.
\end{enumerate}

\end{problem}

\noindent Voilà les questions que de Méré apporta à Pascal en 1654, et la partie (iv) est
\textit{le} problème des partis --- partager équitablement les mises exige de raisonner sur des
parties qui n'ont jamais été jouées, et les lettres de Pascal et Fermat qui le résolvent sont les
actes fondateurs de la théorie des probabilités. Fermat dénombrait ; Pascal récursait ; tous deux
obtinrent la même réponse, et chacune des deux méthodes est devenue un pilier de la
discipline.

\begin{refs}
\item Contexte : Problème des partis --- Wikipédia (\url{https://fr.wikipedia.org/wiki/Probl%C3%A8me_des_partis})
\item Lecture : Joseph K. Blitzstein \& Jessica Hwang, \textit{Introduction to Probability}, ch. 1 (en libre accès (\url{https://probabilitybook.net}))
\item Lecture : William Feller, \textit{An Introduction to Probability Theory and Its Applications}, vol. 1
\end{refs}

\bigskip

\begin{problem}[{Le paradoxe des anniversaires --- Lycée}]
\textbf{PRB-03.} On suppose les anniversaires indépendants et uniformes sur
365 jours.
\begin{enumerate}
\item[(i)] Établissez la probabilité que $ n $ personnes aient toutes des anniversaires distincts :
$ p_n = \prod_{k=0}^{n-1} \big(1 - \frac{k}{365}\big) $, et démontrez que $ n = 23 $ est la
plus petite taille de groupe pour laquelle un anniversaire commun est plus probable
qu'improbable.
\item[(ii)] Expliquez maintenant le phénomène, et pas seulement le nombre : avec $ N $ anniversaires
possibles, utilisez l'inégalité $ 1 - x \le e^{-x} $ (démontrez-la) pour montrer qu'une
collision devient probable dès que $ n \approx 1{,}18\sqrt{N} $, puis utilisez une inégalité
correspondante dans l'autre sens pour montrer que l'ordre $ \sqrt{N} $ est bien le seuil --- la
réponse varie comme la racine carrée de la longueur de l'année, ce qui explique que 23 paraisse
trop petit.
\end{enumerate}

\end{problem}

\noindent Le problème fut posé par Richard von Mises (1939) et popularisé par le manuel
de Feller ; on dit que Harold Davenport le connaissait plus tôt. La loi en $\surd$N de la partie (ii)
en est le contenu véritable : le même calcul régit les collisions dans les tables de hachage et
l'\og{} attaque des anniversaires \fg{} sur les fonctions de hachage cryptographiques, où $ N $ est
astronomiquement grand et où 1,18$\surd$N reste le budget dont un attaquant a besoin.

\begin{refs}
\item Contexte : Paradoxe des anniversaires --- Wikipédia (\url{https://fr.wikipedia.org/wiki/Paradoxe_des_anniversaires})
\item Lecture : Joseph K. Blitzstein \& Jessica Hwang, \textit{Introduction to Probability}, ch. 1 (en libre accès (\url{https://probabilitybook.net}))
\item Cours : MIT OCW 18.600 Probability and Random Variables (\url{https://ocw.mit.edu/courses/18-600-probability-and-random-variables-fall-2019/})
\end{refs}

\bigskip

\begin{problem}[{Monty Hall, démontré --- et non plaidé --- Lycée}]
\textbf{PRB-04.} Une voiture est cachée uniformément au hasard derrière
l'une de trois portes ; les deux autres cachent des chèvres. Vous choisissez la porte 1. Le
présentateur --- qui sait où est la voiture, ouvre toujours une porte que vous n'avez pas choisie,
révèle toujours une chèvre, et choisit uniformément au hasard quand ses deux options cachent des
chèvres --- ouvre la porte 3.
\begin{enumerate}
\item[(i)] Calculez, à partir de la
définition de la probabilité conditionnelle, $ P(\text{voiture derrière la porte 2} \mid
\text{le présentateur ouvre la porte 3}) $, et démontrez que changer de porte fait gagner avec
probabilité $ 2/3 $.
\item[(ii)] Changez maintenant une règle : le présentateur ignore où est la voiture, ouvre uniformément
au hasard l'une des deux portes non choisies, et \textit{se trouve} révéler une chèvre. Démontrez
que la probabilité de gagner en changeant vaut désormais $ 1/2 $.
\item[(iii)] Dites en une phrase ce que ce couple de
calculs démontre au sujet de \og{} la \fg{} réponse à une énigme probabiliste dont le protocole n'est pas
énoncé.
\end{enumerate}

\end{problem}

\noindent Steve Selvin posa le problème dans \textit{The American Statistician} (1975) ; il
explosa en 1990 quand Marilyn vos Savant publia la réponse \og{} changez de porte \fg{} dans
\textit{Parade} et reçut des milliers de lettres --- beaucoup de docteurs en mathématiques ---
soutenant qu'elle avait tort. Paul Erd\H{o}s serait resté incrédule jusqu'à ce qu'on lui montre une
simulation. C'est la démonstration canonique que la plausibilité verbale ne vaut rien en
probabilités : seuls le modèle et le calcul tranchent.

\begin{refs}
\item Contexte : Problème de Monty Hall --- Wikipédia (\url{https://fr.wikipedia.org/wiki/Probl%C3%A8me_de_Monty_Hall})
\item Lecture : Joseph K. Blitzstein \& Jessica Hwang, \textit{Introduction to Probability}, ch. 2 (en libre accès (\url{https://probabilitybook.net}))
\item Cours : Harvard Stat 110 sur edX (Blitzstein) (\url{https://www.edx.org/bio/joseph-blitzstein})
\end{refs}

\bigskip

\begin{problem}[{L'erreur du parieur et la longueur des séries --- Lycée}]
\textbf{PRB-05.} On lance indéfiniment une pièce équilibrée, les lancers étant indépendants.
\begin{enumerate}
\item[(i)] Démontrez, à partir de la définition de l'indépendance, qu'après n'importe quelle suite
d'issues observée --- y compris dix \og{} face \fg{} d'affilée --- la probabilité que le lancer suivant donne
\og{} face \fg{} vaut encore $ 1/2 $.
\item[(ii)] Ceux qui fabriquent de fausses suites \og{} aléatoires \fg{} évitent les longues séries, et c'est
ainsi qu'on les démasque : démontrez qu'en 200 lancers, la probabilité d'observer une série d'au
moins 6 \og{} face \fg{} consécutifs dépasse $ 2/5 $. (Découpez 198 lancers en 33 blocs disjoints de
six et majorez la probabilité qu'aucun bloc ne soit entièrement \og{} face \fg{}.)
\item[(iii)] Dites
précisément ce que la loi des grands nombres ne promet \textit{pas} : montrez que l'espérance de
(nombre de \og{} face \fg{} $ - \; n/2 $) vaut $ 0 $ pour tout $ n $, mais que son écart type vaut
$ \sqrt{n}/2 $ --- la fréquence converge, l'écart absolu non ; les déviations sont diluées, jamais
\og{} corrigées \fg{}.
\end{enumerate}

\end{problem}

\noindent Le 18 août 1913, la roulette de Monte-Carlo sortit le noir 26 fois de suite, et
des joueurs se ruinèrent en misant sur le rouge \og{} parce qu'il était dû \fg{}. Les psychologues Tversky
et Kahneman ont mesuré plus tard la même erreur chez des scientifiques (\og{} Belief in the law of
small numbers \fg{}, 1971). La majoration sur la longueur des séries de la partie (ii) est un
véritable outil médico-légal, et la partie (iii) est l'énoncé honnête du théorème que l'erreur
déforme.

\begin{refs}
\item Contexte : Erreur du parieur --- Wikipédia (\url{https://fr.wikipedia.org/wiki/Erreur_du_parieur})
\item Lecture : William Feller, \textit{An Introduction to Probability Theory and Its Applications}, vol. 1, ch. XIII
\item Cours : MIT OCW 18.600 Probability and Random Variables (\url{https://ocw.mit.edu/courses/18-600-probability-and-random-variables-fall-2019/})
\end{refs}

\bigskip

\begin{problem}[{Bayes au cabinet médical --- Lycée}]
\textbf{PRB-06.} Une maladie touche 1 personne sur 1000. Un test la
détecte avec probabilité $ 0{,}99 $ lorsqu'elle est présente (sensibilité) et donne une fausse
alerte avec probabilité $ 0{,}05 $ lorsqu'elle est absente. Une personne tirée au hasard dans la
population a un test positif.
\begin{enumerate}
\item[(i)] Démontrez le théorème de Bayes
$ P(A\mid B) = \frac{P(B\mid A)\,P(A)}{P(B)} $ à partir de la définition de la probabilité
conditionnelle, y compris la formule des probabilités totales dont vous avez besoin au
dénominateur.
\item[(ii)] Calculez $ P(\text{maladie} \mid \text{test positif}) $ exactement, et expliquez avec des
mots --- en utilisant les effectifs dans une population imaginaire de 100 000 personnes ---
pourquoi la réponse est si loin en dessous de $ 0{,}99 $.
\item[(iii)] Recalculez en supposant que la personne a été testée à cause de symptômes qui portent la
probabilité a priori à $ 1/10 $, et énoncez en une phrase la morale sur les lois a priori.
\end{enumerate}

\end{problem}

\noindent L'essai de Thomas Bayes fut publié à titre posthume en 1763 ; Laplace retrouva
la règle indépendamment et l'utilisa dans toutes les sciences. Dans une célèbre enquête de 1978 à
la faculté de médecine de Harvard, la plupart des médecins répondaient à une version de (ii) par
$\approx$ 95 \% --- l'oubli de la fréquence de base, chez ceux-là mêmes qui prescrivent les tests. La
reformulation en \og{} fréquences naturelles \fg{} de Gigerenzer, celle de la partie (ii), est aujourd'hui
enseignée en médecine parce qu'elle rend la bonne réponse aussi évidente que la mauvaise le
paraissait.

\begin{refs}
\item Contexte : Théorème de Bayes --- Wikipédia (\url{https://fr.wikipedia.org/wiki/Th%C3%A9or%C3%A8me_de_Bayes})
\item Contexte : Oubli de la fréquence de base --- Wikipédia (\url{https://fr.wikipedia.org/wiki/Oubli_de_la_fr%C3%A9quence_de_base})
\item Lecture : Joseph K. Blitzstein \& Jessica Hwang, \textit{Introduction to Probability}, ch. 2 (en libre accès (\url{https://probabilitybook.net}))
\item Cours : Harvard Stat 110 sur edX (Blitzstein) (\url{https://www.edx.org/bio/joseph-blitzstein})
\end{refs}

\bigskip

\begin{problem}[{Deux enfants, deux protocoles --- Lycée, short}]
\textbf{PRB-23.} Une famille a deux enfants, chacun étant indépendamment
garçon ou fille avec probabilité $ 1/2 $. Démontrez que si tout ce qu'on vous dit est qu'au moins
l'un des deux est un garçon, la probabilité conditionnelle que les deux soient des garçons vaut
$ 1/3 $, tandis que si vous rencontrez l'un des deux enfants --- choisi uniformément au hasard --- et
que cet enfant est un garçon, la probabilité que les deux soient des garçons vaut $ 1/2 $.
\end{problem}

\noindent Martin Gardner posa l'énigme dans sa chronique de \textit{Scientific American} en
1959, donna la réponse $ 1/3 $, puis concéda que la question telle qu'elle est formulée ne
détermine pas la réponse : ce qui compte n'est pas le fait que vous apprenez, mais la procédure
qui l'a produit. Les deux calculs présentés ici sont la plus petite illustration complète de la
morale de PRB-04 --- en probabilités, une énigme sans protocole énoncé n'est pas encore un
problème.

\begin{refs}
\item Contexte : Paradoxe des deux enfants --- Wikipédia (\url{https://fr.wikipedia.org/wiki/Paradoxe_des_deux_enfants})
\item Lecture : Joseph K. Blitzstein \& Jessica Hwang, \textit{Introduction to Probability}, ch. 2 (en libre accès (\url{https://probabilitybook.net}))
\end{refs}

\bigskip

\begin{problem}[{Des dés qui se battent en rond --- Lycée, short}]
\textbf{PRB-24.} Quatre dés équilibrés à six faces portent les valeurs
$ A = (4,4,4,4,0,0) $, $ B = (3,3,3,3,3,3) $, $ C = (6,6,2,2,2,2) $,
$ D = (5,5,5,1,1,1) $ ; deux joueurs lancent chacun un dé et le plus grand nombre l'emporte.
Démontrez que
\[ P(A \text{ bat } B) = P(B \text{ bat } C) = P(C \text{ bat } D)
 = P(D \text{ bat } A) = \tfrac{2}{3}, \]
de sorte que \og{} a plus de chances de gagner contre \fg{} n'est pas une relation transitive.
\end{problem}

\noindent Ce sont les dés de Bradley Efron, rendus célèbres par Martin Gardner dans
\textit{Scientific American} en 1970. Ils comptent parce que presque tout le monde suppose en
silence que \og{} meilleur que \fg{} doit être transitif --- l'hypothèse qui sous-tend le classement des
joueurs, des produits et des options de vote --- et voici un contre-exemple en deux lignes ; l'astuce
est qu'une probabilité compare deux lois, et que de telles comparaisons ne s'enchaînent pas
nécessairement.

\begin{refs}
\item Contexte : Dés non transitifs --- Wikipédia (\url{https://fr.wikipedia.org/wiki/D%C3%A9s_non_transitifs})
\item Lecture : Martin Gardner, \textit{The Colossal Book of Mathematics} (chapitre sur les paradoxes non transitifs)
\end{refs}

\bigskip

\begin{problem}[{Le problème des chapeaux --- Lycée}]
\textbf{PRB-25.} À une soirée, $ n $ invités déposent leur chapeau au vestiaire, et un
préposé distrait les leur rend dans un ordre uniformément aléatoire --- c'est-à-dire selon une
permutation uniformément aléatoire de $ \{1, 2, \dots, n\} $. On dit qu'un invité est une
\textit{rencontre} s'il reçoit son propre chapeau.
\begin{enumerate}
\item[(i)] Par inclusion-exclusion, démontrez que la probabilité qu'il n'y ait aucune rencontre vaut
\[ p_n = \sum_{k=0}^{n} \frac{(-1)^k}{k!} , \]
de sorte que le nombre de permutations sans point fixe (un \textit{dérangement}) vaut
$ D_n = n!\, p_n $.
\item[(ii)] Déduisez-en que $ p_n \to e^{-1} $ et que $ |p_n - e^{-1}| < \frac{1}{(n+1)!} $ --- la
réponse est donc essentiellement la même pour 10 invités et pour 10 millions.
\item[(iii)] Démontrez la récurrence $ D_n = (n-1)\big( D_{n-1} + D_{n-2} \big) $ pour $ n \ge 3 $ par
un argument de dénombrement, et non par un calcul algébrique sur la formule.
\item[(iv)] Démontrez que la probabilité d'avoir exactement $ m $ rencontres vaut
$ \frac{1}{m!} \sum_{k=0}^{n-m} \frac{(-1)^k}{k!} $, et montrez que lorsque $ n \to \infty $
cette quantité tend vers $ \frac{e^{-1}}{m!} $ --- la loi de Poisson de PRB-28 de paramètre 1, en
accord avec le nombre moyen de rencontres calculé en PRB-07.
\end{enumerate}

\end{problem}

\noindent Pierre Rémond de Montmort publia ceci sous le nom de \textit{problème des
rencontres} dans son \textit{Essay d'analyse sur les jeux de hazard} (1708), où il modélisait un
jeu de cartes appelé le treize ; Nicolas Bernoulli puis Euler y revinrent, et la somme alternée
de (i) est l'une des premières apparitions imprimées du principe d'inclusion-exclusion. La chute
de (ii) fait la célébrité du problème : une probabilité qui refuse obstinément de dépendre de la
taille de la soirée, fixée à $ 1/e $ dès $ n = 6 $ environ et jusqu'aux limites de la
mesure.

\begin{refs}
\item Contexte : Problème des rencontres --- Wikipédia (\url{https://fr.wikipedia.org/wiki/Probl%C3%A8me_des_rencontres})
\item Contexte : Principe d'inclusion-exclusion --- Wikipédia (\url{https://fr.wikipedia.org/wiki/Principe_d'inclusion-exclusion})
\item Lecture : William Feller, \textit{An Introduction to Probability Theory and Its Applications}, vol. 1, ch. IV
\item Cours : MIT OCW 18.600 Probability and Random Variables (\url{https://ocw.mit.edu/courses/18-600-probability-and-random-variables-fall-2019/})
\end{refs}

\bigskip

\begin{problem}[{Ce que trois axiomes imposent déjà --- Lycée, short}]
\textbf{PRB-35.} Les axiomes de Kolmogorov disent seulement ceci : tout
événement $ A $ reçoit un nombre $ P(A) \ge 0 $, l'événement certain vérifie
$ P(\Omega) = 1 $, et $ P(A_1 \cup A_2 \cup \cdots) = \sum_{n} P(A_n) $ dès que les événements
$ A_n $ sont deux à deux disjoints. En n'utilisant rien d'autre, démontrez la règle du
complémentaire $ P(A^c) = 1 - P(A) $, la croissance ($ A \subseteq B $ entraîne
$ P(A) \le P(B) $), la formule à deux événements
$ P(A \cup B) = P(A) + P(B) - P(A \cap B) $ et l'inégalité de Boole
$ P\big( \bigcup_n A_n \big) \le \sum_n P(A_n) $ ; démontrez ensuite qu'aucune probabilité
n'attribue le \textit{même} nombre à chacun des événements $ \{1\}, \{2\}, \{3\}, \dots $ sur
l'univers $ \Omega = \{1,2,3,\dots\} $, de sorte que \og{} tirer un entier naturel au hasard \fg{} ne
décrit rien du tout.
\end{problem}

\noindent Andreï Kolmogorov publia ces axiomes dans \textit{Grundbegriffe der
Wahrscheinlichkeitsrechnung} (1933), réglant pour les probabilités la part du sixième problème
de Hilbert (1900) qui réclamait un traitement axiomatique de la discipline ; en une décennie, tout
résultat de ce domaine devint un théorème de théorie de la mesure. La dernière partie du problème
montre que le troisième axiome --- l'additivité sur une infinité \textit{dénombrable} d'événements
disjoints --- est un engagement réel et non une commodité comptable : c'est exactement lui qui tue
la loi uniforme sur les entiers, et c'est pourquoi Bruno de Finetti (PRB-18) plaidait pour une
théorie seulement finiment additive. Le même axiome impose une restriction de plus, invisible à ce
niveau : aucune probabilité dénombrablement additive et invariante par translation ne peut être
définie sur \textit{toutes} les parties du cercle, comme Giuseppe Vitali l'a montré en 1905, ce qui
explique que les axiomes parlent d'une famille d'événements choisie plutôt que de tous les
ensembles sans exception.

\begin{refs}
\item Contexte : Axiomes des probabilités --- Wikipédia (\url{https://fr.wikipedia.org/wiki/Axiomes_des_probabilit%C3%A9s})
\item Contexte : Ensemble de Vitali --- Wikipédia (\url{https://fr.wikipedia.org/wiki/Ensemble_de_Vitali})
\item Lecture : A. N. Kolmogorov, \textit{Foundations of the Theory of Probability} (1933 ; traduction anglaise, Chelsea, 1956)
\end{refs}

\bigskip

\begin{problem}[{Le problème du scrutin --- Lycée}]
\textbf{PRB-36.} Lors d'une élection, le candidat A recueille $ p $ voix et le candidat B en
recueille $ q $, avec $ p > q $. Les $ p + q $ bulletins sont dépouillés un à un dans un
ordre uniformément aléatoire, les $ \binom{p+q}{p} $ ordres étant équiprobables. On représente un
dépouillement par un chemin qui monte à chaque voix pour A et descend à chaque voix pour B.
\begin{enumerate}
\item[(a)] Démontrez le théorème du scrutin : la probabilité que A soit strictement en tête de B à chaque
étape du dépouillement vaut
\[ \frac{p-q}{p+q} . \]
(Réfléchissez le segment initial de chaque mauvais chemin par rapport à l'axe horizontal jusqu'au
premier instant d'égalité entre les deux candidats, et montrez que cela apparie bijectivement les
mauvais chemins commençant par une voix pour A avec ceux commençant par une voix pour B.)
\item[(b)] Démontrez-le une seconde fois par le lemme du cycle : écrivez les $ p+q $ bulletins autour
d'un cercle et montrez que, parmi les $ p+q $ rotations d'une disposition fixée, exactement
$ p - q $ donnent un dépouillement où A est en tête du début à la fin.
\item[(c)] Démontrez la forme faible --- le nombre d'ordres dans lesquels A n'est jamais \textit{derrière}
vaut $ \frac{p+1-q}{p+1}\binom{p+q}{q} $ --- et déduisez-en que pour $ p = q = n $ le nombre de
tels ordres est le nombre de Catalan $ \frac{1}{n+1}\binom{2n}{n} $.
\item[(d)] Remarquez ce dont la réponse ne dépend pas. Calculez la probabilité que le vainqueur soit en
tête du premier au dernier bulletin dans un électorat d'un million de personnes se partageant
600 000 contre 400 000, et comparez avec un village de dix électeurs se partageant 6 contre
4 ; montrez ensuite qu'une marge d'une seule voix ($ p = q+1 $) rend une avance de bout en bout
essentiellement impossible, avec probabilité $ 1/(p+q) $.
\end{enumerate}

\end{problem}

\noindent Joseph Bertrand posa la question dans les \textit{Comptes rendus} en 1887 et la
résolut par une récurrence, en remarquant que la simplicité de la réponse appelait une
démonstration directe ; W. A. Whitworth avait en fait publié le résultat dès 1878. Quelques
semaines après Bertrand, Désiré André fournit un argument direct, et tous les manuels depuis
l'appellent \og{} le principe de réflexion \fg{} --- mais Marc Renault, lisant l'original français en 2008, a
montré qu'André n'avait rien réfléchi : sa bijection permute les voix, et la démonstration par
réflexion n'entre dans les archives qu'en 1923. Le lemme du cycle de la partie (b) est dû à
Dvoretzky et Motzkin (1947), et le théorème lui-même est le germe de la théorie des fluctuations
de PRB-38 et des apparitions des nombres de Catalan dans toute la combinatoire.

\begin{refs}
\item Contexte : Problème du scrutin --- Wikipédia (\url{https://fr.wikipedia.org/wiki/Probl%C3%A8me_du_scrutin})
\item Contexte : Nombre de Catalan --- Wikipédia (\url{https://fr.wikipedia.org/wiki/Nombre_de_Catalan})
\item Article : M. Renault, \og{} Lost (and found) in translation: André's actual method and its application to the generalized ballot problem \fg{}, Amer. Math. Monthly 115 (2008), 358--363
\item Lecture : William Feller, \textit{An Introduction to Probability Theory and Its Applications}, vol. 1, ch. III
\item Voir aussi : PUZ-51 (\url{puzzles.html#PUZ-51}) --- la page Casse-tête, où le même dénombrement est posé comme une énigme.
\end{refs}

\bigskip


\section*{Licence}

\begin{problem}[{Linéarité de l'espérance --- Licence}]
\textbf{PRB-07.} \begin{enumerate}
\item[(i)] Démontrez que $ E[X+Y] = E[X] + E[Y] $ pour deux variables aléatoires quelconques
d'espérance finie --- sans aucune hypothèse d'indépendance --- et expliquez pourquoi c'est remarquable
lorsque $ X $ et $ Y $ sont fortement dépendantes. Utilisez ensuite des variables indicatrices
pour démontrer :
\item[(ii)] le théorème du collectionneur de vignettes --- en tirant uniformément avec remise parmi $ n $
types de vignettes, le nombre moyen de tirages nécessaires pour les collecter toutes vaut
$ n H_n = n\big(1 + \frac12 + \cdots + \frac1n\big) \approx n\ln n $ ;
\item[(iii)] une permutation uniformément aléatoire de $ n $ lettres a exactement $ 1 $ point fixe en
moyenne, et cela pour tout $ n $ ;
\item[(iv)] défi : son nombre moyen de cycles vaut $ H_n $.
\end{enumerate}

\end{problem}

\noindent Le problème des rencontres de Montmort (1708) --- combien d'invités récupèrent
leur propre chapeau ? --- c'est la partie (iii) en costume d'époque. La linéarité de l'espérance est
la plus puissante des astuces bon marché des probabilités : des sommes d'indicatrices dépendantes
qu'il serait sans espoir de traiter conjointement deviennent triviales une par une, et Erd\H{o}s a
bâti la méthode probabiliste exactement sur ce geste.

\begin{refs}
\item Contexte : Espérance mathématique --- Wikipédia (\url{https://fr.wikipedia.org/wiki/Esp%C3%A9rance_math%C3%A9matique})
\item Contexte : Problème du collectionneur de vignettes --- Wikipédia (\url{https://fr.wikipedia.org/wiki/Probl%C3%A8me_du_collectionneur_de_vignettes})
\item Lecture : Joseph K. Blitzstein \& Jessica Hwang, \textit{Introduction to Probability}, ch. 4 (en libre accès (\url{https://probabilitybook.net}))
\item Cours : MIT OCW 18.600 Probability and Random Variables (\url{https://ocw.mit.edu/courses/18-600-probability-and-random-variables-fall-2019/})
\end{refs}

\bigskip

\begin{problem}[{Tchebychev et la loi faible --- Licence}]
\textbf{PRB-08.} \begin{enumerate}
\item[(i)] Démontrez l'inégalité de Markov : pour $ X\ge 0 $
et $ a>0 $, $ P(X\ge a) \le E[X]/a $.
\item[(ii)] Déduisez-en l'inégalité de Bienaymé-Tchebychev
$ P(|X-\mu| \ge k\sigma) \le 1/k^2 $.
\item[(iii)] Démontrez que la variance d'une somme de variables aléatoires deux à deux indépendantes est la
somme des variances, et déduisez-en la loi faible des grands nombres pour des variables i.i.d. de
variance finie : $ P\big( \big| \frac{S_n}{n} - \mu \big| \ge \varepsilon \big) \to 0 $.
\item[(iv)] Appliquez-la : combien de lancers d'une pièce équilibrée Tchebychev exige-t-il pour garantir,
avec probabilité au moins $ 0{,}95 $, que la fréquence observée de \og{} face \fg{} soit à $ 0{,}01 $
près de $ 1/2 $ ? Comparez avec le nombre bien plus petit que suggère le théorème central limite
(PRB-14), et expliquez l'écart.
\end{enumerate}

\end{problem}

\noindent Jacques Bernoulli démontra la première loi des grands nombres dans
\textit{Ars Conjectandi} (publié en 1713) après vingt ans d'efforts --- il l'appelait son \og{} théorème
d'or \fg{}. L'inégalité de Tchebychev de 1867 la réduisit à quatre lignes. Ce couple d'inégalités est
l'atome d'hydrogène de la concentration de la mesure : grossières, universelles, et le patron de
toutes les majorations plus fines depuis.

\begin{refs}
\item Contexte : Inégalité de Bienaymé-Tchebychev --- Wikipédia (\url{https://fr.wikipedia.org/wiki/In%C3%A9galit%C3%A9_de_Bienaym%C3%A9-Tchebychev})
\item Contexte : Loi des grands nombres --- Wikipédia (\url{https://fr.wikipedia.org/wiki/Loi_des_grands_nombres})
\item Lecture : Joseph K. Blitzstein \& Jessica Hwang, \textit{Introduction to Probability} (en libre accès (\url{https://probabilitybook.net}))
\item Cours : MIT OCW 18.600 Probability and Random Variables (\url{https://ocw.mit.edu/courses/18-600-probability-and-random-variables-fall-2019/})
\end{refs}

\bigskip

\begin{problem}[{Attendre le premier succès --- Licence}]
\textbf{PRB-09.} Des épreuves indépendantes réussissent avec probabilité $ p $.
\begin{enumerate}
\item[(i)] Établissez la loi géométrique $ P(T = k) = (1-p)^{k-1}p $ de l'instant du premier succès et
vérifiez qu'elle somme à 1.
\item[(ii)] Démontrez $ E[T] = 1/p $ de deux façons : en sommant la série, et par l'identité de
conditionnement $ E[T] = 1 + (1-p)E[T] $ --- et justifiez rigoureusement pourquoi cette identité
est vraie et pourquoi il est légitime de la résoudre.
\item[(iii)] Démontrez que la loi géométrique est sans mémoire ---
$ P(T > m+n \mid T > m) = P(T > n) $ --- et démontrez que c'est la \textit{seule} loi sur
$ \{1,2,3,\dots\} $ ayant cette propriété.
\item[(iv)] Établissez la loi binomiale négative de l'instant du $ r $-ième succès, présentez-la comme
une somme de $ r $ variables géométriques indépendantes, et calculez son espérance et sa
variance.
\end{enumerate}

\end{problem}

\noindent La loi binomiale négative apparaît chez Pascal --- on l'appelle encore la loi de
Pascal --- et la caractérisation par l'absence de mémoire de (iii) est l'ombre discrète de la
propriété qui définit la loi exponentielle, la raison pour laquelle les temps d'attente
géométriques sont les atomes de la théorie des chaînes de Markov. Les deux calculs de $ E[T] $
en (ii) --- force brute contre auto-similarité --- sont une première leçon d'argument de
renouvellement.

\begin{refs}
\item Contexte : Loi géométrique --- Wikipédia (\url{https://fr.wikipedia.org/wiki/Loi_g%C3%A9om%C3%A9trique})
\item Contexte : Loi binomiale négative --- Wikipédia (\url{https://fr.wikipedia.org/wiki/Loi_binomiale_n%C3%A9gative})
\item Lecture : Joseph K. Blitzstein \& Jessica Hwang, \textit{Introduction to Probability}, ch. 4 (en libre accès (\url{https://probabilitybook.net}))
\end{refs}

\bigskip

\begin{problem}[{La ruine du joueur --- Licence}]
\textbf{PRB-10.} Un joueur part avec $ i $ unités et mise une unité par
manche, gagnant avec probabilité $ p $ et perdant avec probabilité $ q = 1-p $, de façon
indépendante ; la partie s'arrête à la fortune $ 0 $ (ruine) ou $ N $ (victoire).
\begin{enumerate}
\item[(i)] Écrivez l'équation aux différences vérifiée par la probabilité de ruine $ r_i $,
résolvez-la pour $ p = 1/2 $ puis pour $ p \ne 1/2 $, et démontrez que votre solution est
l'unique solution satisfaisant les conditions aux limites.
\item[(ii)] Calculez la durée moyenne de la partie.
\item[(iii)] Démontrez que contre un adversaire infiniment riche ($ N \to \infty $) à jeu équitable, la
ruine est certaine.
\item[(iv)] Chiffrez l'avantage de la maison : avec $ p = 18/38 $ (roulette américaine, mise sur le
rouge), comparez les probabilités de ruine d'un joueur qui possède 50 unités et veut atteindre 100
par mises unitaires avec celles d'un joueur qui mise ses 50 unités sur un seul tour --- et démontrez
laquelle des deux stratégies est la meilleure.
\end{enumerate}

\end{problem}

\noindent Pascal posa le problème de la ruine à Fermat en 1656, et Huygens l'inclut dans
le premier traité imprimé de probabilités (1657) ; la résolution par équation aux différences est
essentiellement celle d'Abraham de Moivre. La partie (iv) --- le jeu prudent perd plus sûrement que
le jeu audacieux face à un jeu défavorable --- est un théorème que les joueurs paient pour
redécouvrir, et l'ensemble du problème est la porte d'entrée des chaînes de Markov absorbantes et
des idées d'arrêt optionnel de PRB-16.

\begin{refs}
\item Contexte : Ruine du joueur --- Wikipédia (\url{https://fr.wikipedia.org/wiki/Ruine_du_joueur})
\item Lecture : William Feller, \textit{An Introduction to Probability Theory and Its Applications}, vol. 1, ch. XIV
\item Lecture : Geoffrey Grimmett \& David Stirzaker, \textit{Probability and Random Processes}
\end{refs}

\bigskip

\begin{problem}[{Le paradoxe de Saint-Pétersbourg --- Licence}]
\textbf{PRB-11.} On lance une pièce équilibrée jusqu'au premier \og{} face \fg{} ;
si cela se produit au lancer $ k $, le casino verse $ 2^k $ ducats.
\begin{enumerate}
\item[(i)] Démontrez que le gain espéré est infini --- et pourtant personne ne paierait 1000 ducats pour
jouer une fois.
\item[(ii)] La résolution de Daniel Bernoulli : un joueur de fortune $ w $ évalue l'argent par
$ \ln w $. Montrez que le \textit{gain d'utilité} espéré du jeu est fini, et calculez
(numériquement) le droit d'entrée équitable qu'il implique pour un joueur de fortune 1000.
\item[(iii)] La résolution de Feller : l'espérance est simplement le mauvais résumé pour un jeu unique à
queue lourde. Rendez son énoncé précis --- pour $ n $ parties indépendantes de gain total
$ S_n $, le droit d'entrée \og{} équitable \fg{} par partie croît comme $ \log_2 n $ --- et démontrez la
moitié accessible : si chaque partie coûte un montant fixe $ c $, le casino gagne à la longue,
aussi grand que $ c $ soit choisi.
\item[(iv)] Expliquez pourquoi (i)--(iii) ne se contredisent pas.
\end{enumerate}

\end{problem}

\noindent Nicolas Bernoulli proposa le jeu dans une lettre à Montmort en 1713 ; son cousin
Daniel publia la résolution par l'utilité dans le journal de l'Académie de Saint-Pétersbourg
(1738), fondant la théorie de l'utilité espérée et, avec elle, une grande part de l'économie.
L'analyse de Feller au XXe siècle a recadré la question : le paradoxe ne porte pas sur
la psychologie mais sur ce que la loi des grands nombres autorise réellement quand la variance est
infinie.

\begin{refs}
\item Contexte : Paradoxe de Saint-Pétersbourg --- Wikipédia (\url{https://fr.wikipedia.org/wiki/Paradoxe_de_Saint-P%C3%A9tersbourg})
\item Lecture : William Feller, \textit{An Introduction to Probability Theory and Its Applications}, vol. 1
\item Cours : MIT OCW 18.600 Probability and Random Variables (\url{https://ocw.mit.edu/courses/18-600-probability-and-random-variables-fall-2019/})
\end{refs}

\bigskip

\begin{problem}[{L'aiguille de Buffon --- Licence}]
\textbf{PRB-12.} On laisse tomber une aiguille de longueur $ \ell $ sur un
parquet rayé de droites parallèles distantes de $ d \ge \ell $. (i) Construisez le modèle --- la
distance du milieu de l'aiguille à la droite la plus proche uniforme sur $ [0, d/2] $, son angle
uniforme sur $ [0, \pi/2] $, les deux indépendants --- et démontrez
\[ P(\text{l'aiguille coupe une droite}) = \frac{2\ell}{\pi d}. \]
(ii) Transformez cela en un estimateur de $\pi$ à partir de $ n $ chutes indépendantes ; calculez la
variance de l'estimateur et utilisez Tchebychev (PRB-08) pour majorer le nombre de chutes
garantissant deux décimales correctes avec probabilité $ 0{,}99 $. (iii) La \og{} nouille de
Buffon \fg{} de Barbier : démontrez que pour une courbe rigide de forme quelconque et de longueur
$ L $, le \textit{nombre moyen} d'intersections avec les droites vaut $ \frac{2L}{\pi d} $ --- en
n'utilisant rien d'autre que la linéarité de l'espérance (PRB-07).
\end{problem}

\noindent L'\textit{Essai d'arithmétique morale} de Buffon (1777) a fait de ceci le premier
problème de probabilités géométriques. En 1901, Lazzarini annonça avoir estimé $\pi$ à 355/113 à
partir de 3408 chutes --- une précision si improbable que son expérience est aujourd'hui l'exemple
scolaire de données trop belles pour être vraies ; votre majoration de la variance en (ii) est
l'outil qui le confond. L'argument de Barbier en 1860, celui de la partie (iii), est un joyau de
la méthode des indicatrices.

\begin{refs}
\item Contexte : Aiguille de Buffon --- Wikipédia (\url{https://fr.wikipedia.org/wiki/Aiguille_de_Buffon})
\item Contexte : La nouille de Buffon --- Wikipédia (en anglais) (\url{https://en.wikipedia.org/wiki/Buffon%27s_noodle})
\item Lecture : Daniel A. Klain \& Gian-Carlo Rota, \textit{Introduction to Geometric Probability}
\end{refs}

\bigskip

\begin{problem}[{L'ivrogne de Pólya : marches aléatoires en dimension d --- Licence}]
\textbf{PRB-13.} La marche aléatoire simple sur $ \mathbb{Z}^d $ saute
vers un plus proche voisin choisi uniformément, indépendamment et indéfiniment ; on dit que la
marche est \textit{récurrente} si elle revient à son point de départ avec probabilité 1.
\begin{enumerate}
\item[(i)] Démontrez le critère : la marche est récurrente si et seulement si
$ \sum_n P(S_{2n} = 0) = \infty $.
\item[(ii)] Pour $ d=1 $, montrez que
$ P(S_{2n}=0) = \binom{2n}{n}2^{-2n} \sim \frac{1}{\sqrt{\pi n}} $ (Stirling, ou PRB-14), et
concluez à la récurrence.
\item[(iii)] Pour $ d=2 $, démontrez que $ P(S_{2n}=0) \sim \frac{c}{n} $ (une identité combinatoire
ramène le calcul à deux marches unidimensionnelles indépendantes) --- récurrence de nouveau.
\item[(iv)] Pour $ d=3 $, montrez que $ P(S_{2n}=0) = O(n^{-3/2}) $ et concluez à la transience : avec
probabilité strictement positive, la marche ne rentre jamais chez elle.
\end{enumerate}

\end{problem}

\noindent George Pólya démontra ceci en 1921, inspiré --- disait-il --- par le fait de croiser
sans cesse le même couple en se promenant dans les bois près de Zurich. Le résumé de Kakutani est
immortel : \og{} Un homme ivre finira par retrouver son chemin, mais un oiseau ivre risque de se
perdre à jamais. \fg{} Le seuil de dimension entre récurrence et transience est la première apparition
d'une transition de phase dans cette discipline --- un thème qui revient en pleine force en PRB-17
et PRB-20.

\begin{refs}
\item Contexte : Marche aléatoire --- Wikipédia (\url{https://fr.wikipedia.org/wiki/Marche_al%C3%A9atoire})
\item Article : G. Pólya, \og{} Über eine Aufgabe der Wahrscheinlichkeitsrechnung betreffend die Irrfahrt im Straßennetz \fg{}, Math. Ann. 84 (1921)
\item Lecture : Peter G. Doyle \& J. Laurie Snell, \textit{Random Walks and Electric Networks} (en libre accès)
\item Lecture : Rick Durrett, \textit{Probability: Theory and Examples} (PDF gratuit, Duke (\url{https://sites.math.duke.edu/~rtd/PTE/PTE5_011119.pdf}))
\end{refs}

\bigskip

\begin{problem}[{La courbe en cloche gagne sa place : de Moivre-Laplace --- Licence}]
\textbf{PRB-14.} Soit $ S_n $ le nombre de \og{} face \fg{} en $ n $ lancers
d'une pièce équilibrée.
\begin{enumerate}
\item[(i)] Démontrez une formule de Stirling faible $ n! \sim C\sqrt{n}\,(n/e)^n $ en comparant
$ \ln n! $ à une intégrale, puis identifiez $ C = \sqrt{2\pi} $ grâce au produit de Wallis (ou
à l'intégrale de Gauss, CAL-15 sur la page
Calcul différentiel et intégral (\url{applied-calculus.html})).
\item[(ii)] Démontrez le théorème limite local : pour $ k = n/2 + x\sqrt{n}/2 $ avec $ x $ borné,
\[ P(S_n = k) \sim \frac{2}{\sqrt{2\pi n}}\; e^{-x^2/2} . \]
\item[(iii)] Sommez : déduisez-en le théorème de Moivre-Laplace
$ P\big( a \le \frac{S_n - n/2}{\sqrt{n}/2} \le b \big) \to
\frac{1}{\sqrt{2\pi}}\int_a^b e^{-x^2/2}\,dx $.
\item[(iv)] Énoncez soigneusement --- sans démonstration ---
le théorème central limite général pour les suites i.i.d. de variance finie, et expliquez ce que
signifie ici l'\og{} universalité \fg{} : quels détails de la pièce votre limite a-t-elle oubliés ?
\end{enumerate}

\end{problem}

\noindent De Moivre publia l'approximation en 1733 sous forme de supplément privé, plus
tard intégré à \textit{The Doctrine of Chances} ; Laplace l'étendit et fit de la courbe normale le
cheval de bataille de la statistique céleste et sociale. Le TCL général --- Lindeberg, Lévy, années
1920 --- se démontre par les fonctions caractéristiques (voir Durrett), mais le cas de la pièce que
vous traitez ici contient le phénomène tout entier : les détails microscopiques s'effacent, et
$ e^{-x^2/2} $ est ce qui reste.

\begin{refs}
\item Contexte : Théorème de Moivre-Laplace --- Wikipédia (\url{https://fr.wikipedia.org/wiki/Th%C3%A9or%C3%A8me_de_Moivre-Laplace})
\item Contexte : Théorème central limite --- Wikipédia (\url{https://fr.wikipedia.org/wiki/Th%C3%A9or%C3%A8me_central_limite})
\item Lecture : Rick Durrett, \textit{Probability: Theory and Examples}, ch. 3 (PDF gratuit, Duke (\url{https://sites.math.duke.edu/~rtd/PTE/PTE5_011119.pdf}))
\item Cours : MIT OCW 18.600 Probability and Random Variables (\url{https://ocw.mit.edu/courses/18-600-probability-and-random-variables-fall-2019/})
\end{refs}

\bigskip

\begin{problem}[{Seule l'exponentielle oublie --- Licence, short}]
\textbf{PRB-26.} Soit $ T $ une variable aléatoire telle que
$ P(T < \infty) = 1 $ et $ P(T > t) > 0 $ pour tout $ t \ge 0 $, vérifiant la
propriété d'absence de mémoire
\[ P(T > s+t) = P(T > s)\, P(T > t) \qquad \text{pour tous } s,t \ge 0 . \]
Démontrez qu'il existe $ \lambda > 0 $ tel que $ P(T > t) = e^{-\lambda t} $, en prenant
soin de justifier le passage de $ t $ rationnel à $ t $ réel en n'utilisant que la monotonie de
la fonction de survie.
\end{problem}

\noindent C'est le jumeau continu de la caractérisation géométrique de PRB-09(iii), et
c'est la raison pour laquelle la loi exponentielle est incontournable : tout temps d'attente qui ne
vieillit pas est exponentiel. L'équation fonctionnelle $ f(s+t) = f(s)f(t) $ est celle de Cauchy
(1821), et elle possède des solutions monstrueuses non mesurables si l'on n'impose aucune
régularité --- Hamel en a exhibé en 1905 --- de sorte que la monotonie que vous exploitez n'est pas un
détail technique mais toute l'hypothèse. Tout ce qui est bâti sur des temps de séjour exponentiels,
des processus de Poisson aux chaînes de Markov à temps continu et à la théorie des files
d'attente, repose sur ce seul théorème.

\begin{refs}
\item Contexte : Perte de mémoire (probabilités) --- Wikipédia (\url{https://fr.wikipedia.org/wiki/Perte_de_m%C3%A9moire_(probabilit%C3%A9s)})
\item Contexte : Loi exponentielle --- Wikipédia (\url{https://fr.wikipedia.org/wiki/Loi_exponentielle})
\item Lecture : Geoffrey Grimmett \& David Stirzaker, \textit{Probability and Random Processes}
\end{refs}

\bigskip

\begin{problem}[{Où tombe le $ k $-ième plus petit point --- Licence, short}]
\textbf{PRB-27.} Soient $ U_1, \dots, U_n $ indépendantes et uniformes sur
$ (0,1) $, et soient $ U_{(1)} < U_{(2)} < \cdots < U_{(n)} $ ces mêmes nombres rangés
par ordre croissant. Démontrez que $ U_{(k)} $ admet la densité
\[ f_{(k)}(x) = \frac{n!}{(k-1)!\,(n-k)!}\; x^{k-1}(1-x)^{n-k}, \qquad 0 < x < 1, \]
et déduisez-en que $ E\big[U_{(k)}\big] = \frac{k}{n+1} $ --- les $ n $ points découpent
l'intervalle en $ n+1 $ morceaux de même longueur moyenne.
\end{problem}

\noindent Ce sont les statistiques d'ordre, et la densité ci-dessus est la loi bêta que
Bayes intégrait déjà en 1763. Le fait que les espacements aient la même longueur moyenne est
l'ancêtre continu de l'estimateur du char d'assaut allemand (STA-11 sur la page
Statistique (\url{applied-statistics.html})), où la même correction \og{} ajoutez l'écart moyen \fg{}
transforme le maximum de l'échantillon en estimateur sans biais ; il sous-tend aussi les positions
de tracé d'un diagramme quantile-quantile et la statistique de Kolmogorov-Smirnov.

\begin{refs}
\item Contexte : Statistique d'ordre --- Wikipédia (\url{https://fr.wikipedia.org/wiki/Statistique_d'ordre})
\item Lecture : H. A. David \& H. N. Nagaraja, \textit{Order Statistics}, 3e éd., ch. 1--2
\item Lecture : Joseph K. Blitzstein \& Jessica Hwang, \textit{Introduction to Probability}, ch. 8 (en libre accès (\url{https://probabilitybook.net}))
\end{refs}

\bigskip

\begin{problem}[{La loi des petits nombres --- Licence}]
\textbf{PRB-28.} Soient $ X_1, \dots, X_n $ indépendantes avec $ X_i \sim
\mathrm{Bernoulli}(p_i) $, soit $ S_n = X_1 + \cdots + X_n $, et posons
$ \lambda = p_1 + \cdots + p_n $.
\begin{enumerate}
\item[(i)] Démontrez le théorème limite de Poisson sous sa forme classique : si $ p_i = \lambda/n $
pour tout $ i $, alors pour chaque $ k $ fixé
\[ P(S_n = k) \longrightarrow \frac{e^{-\lambda}\lambda^k}{k!} \qquad (n \to \infty) . \]
\item[(ii)] Démontrez qu'une variable de Poisson($ \lambda $) a une espérance et une variance toutes
deux égales à $ \lambda $, calculez sa fonction génératrice des probabilités, et déduisez-en
qu'une somme de variables de Poisson indépendantes suit une loi de Poisson de paramètre la somme
des paramètres.
\item[(iii)] Renforcez la limite en une inégalité --- le théorème de Le Cam : la distance en variation totale
entre la loi de $ S_n $ et Poisson($ \lambda $) vérifie
\[ \sum_{k \ge 0} \big| P(S_n = k) - e^{-\lambda}\lambda^k/k! \big| \;\le\; 2\sum_{i=1}^{n} p_i^2 . \]
(Couplez chaque $ X_i $ avec une variable de Poisson($ p_i $) sur un même espace probabilisé de
sorte qu'elles diffèrent avec probabilité au plus $ p_i^2 $, puis comparez les sommes.) Remarquez
qu'aucun \og{} $ n \to \infty $ \fg{} n'apparaît : la majoration est bonne dès que les probabilités
individuelles sont petites, si inégales soient-elles.
\item[(iv)] Testez la théorie sur le jeu de données fondateur. Ladislaus von Bortkiewicz a relevé, pour 10
corps d'armée prussiens sur les 20 années 1875--1894, le nombre de soldats tués par ruade de
cheval : 109 années-corps avec 0 mort, 65 avec 1, 22 avec 2, 3 avec 3, 1 avec 4. Estimez
$ \lambda $ à partir du total, comparez les 200 effectifs attendus sous la loi de Poisson aux
effectifs observés, et expliquez précisément quelles hypothèses de (iii) rendent le modèle
plausible ici.
\end{enumerate}

\end{problem}

\noindent Siméon Denis Poisson établit la loi en 1837 dans une étude sur la probabilité
des condamnations par les jurys français --- une application dont personne ne se souvient, dans un
livre que presque personne n'a lu à l'époque. Elle ne devint célèbre qu'après que
\textit{Das Gesetz der kleinen Zahlen} de Bortkiewicz (1898) l'eut ajustée aux données des ruades de
cheval et eut donné son nom au sujet : des événements rares, beaucoup d'occasions indépendantes,
une moyenne stable. La majoration de Le Cam de 1960, celle de la partie (iii), est la forme moderne
du théorème --- elle remplace une limite par une estimation d'erreur, et c'est elle qui permet
d'appliquer l'approximation de Poisson à une collection finie, hétérogène et désordonnée de petits
risques.

\begin{refs}
\item Contexte : Théorème limite de Poisson --- Wikipédia (en anglais) (\url{https://en.wikipedia.org/wiki/Poisson_limit_theorem})
\item Contexte : Inégalité de Le Cam --- Wikipédia (\url{https://fr.wikipedia.org/wiki/In%C3%A9galit%C3%A9_de_Le_Cam})
\item Article : L. Le Cam, \og{} An approximation theorem for the Poisson binomial distribution \fg{}, Pacific J. Math. 10 (1960), 1181--1197
\item Lecture : William Feller, \textit{An Introduction to Probability Theory and Its Applications}, vol. 1, ch. VI
\end{refs}

\bigskip

\begin{problem}[{La moyenne qui n'apprend rien --- Licence, short}]
\textbf{PRB-37.} Soient $ X_1, X_2, \dots, X_n $ indépendantes de densité de
Cauchy standard $ f(x) = \frac{1}{\pi(1+x^2)} $ sur $ \mathbb{R} $. Démontrez que la moyenne
empirique $ \bar{X}_n = (X_1 + \cdots + X_n)/n $ suit \textit{exactement} la loi de Cauchy standard
pour tout $ n $ --- par exemple en calculant la fonction caractéristique
$ E\big[e^{itX_1}\big] = e^{-|t|} $ --- de sorte que moyenner un million d'observations ne vaut pas
mieux que garder la première, et identifiez précisément quelle hypothèse de la loi des grands
nombres (PRB-08) a été violée.
\end{problem}

\noindent Poisson écrivit cette densité en 1824 comme une curiosité ; elle porte le nom de
Cauchy à cause de son échange de 1853 avec Irénée-Jules Bienaymé au sujet de la méthode des
moindres carrés, où cet exemple précis servit de munition --- Bienaymé défendait les moindres carrés,
Cauchy produisait des erreurs pour lesquelles les moyennes empiriques ne se stabilisent pas. La loi
de Cauchy est la loi stable symétrique d'indice $ \alpha = 1 $, et les lois stables sont
exactement les limites possibles de sommes normalisées de variables i.i.d., le théorème central
limite de PRB-14 étant le cas $ \alpha = 2 $ ; tout le reste de la famille a une variance infinie
et une queue qui décroît comme une puissance. Benoit Mandelbrot soutint en 1963 que les cours du
coton se comportent ainsi, et le débat sur les queues lourdes en finance et en assurance n'a jamais
cessé depuis.

\begin{refs}
\item Contexte : Loi de Cauchy (probabilités) --- Wikipédia (\url{https://fr.wikipedia.org/wiki/Loi_de_Cauchy_(probabilit%C3%A9s)})
\item Contexte : Loi stable --- Wikipédia (\url{https://fr.wikipedia.org/wiki/Loi_stable})
\item Lecture : William Feller, \textit{An Introduction to Probability Theory and Its Applications}, vol. 2, ch. VI et XVII
\end{refs}

\bigskip

\begin{problem}[{Réflexion, premier passage et loi de l'arcsinus --- Licence}]
\textbf{PRB-38.} Soient $ S_0 = 0 $ et $ S_k = X_1 + \cdots + X_k $, où les $ X_i $ sont
indépendantes et valent $ +1 $ ou $ -1 $ avec probabilité $ 1/2 $ chacune. Notons
$ M_n = \max_{0 \le k \le n} S_k $ et $ u_{2m} = P(S_{2m} = 0) = \binom{2m}{m} 2^{-2m} $.
\begin{enumerate}
\item[(a)] Démontrez le principe de réflexion : pour des entiers $ a \ge 1 $ et $ x \le a $, les
chemins de $ (0,0) $ à $ (n,x) $ qui touchent le niveau $ a $ sont en bijection avec
\textit{tous} les chemins de $ (0,0) $ à $ (n, 2a - x) $. Déduisez-en
\[ P(M_n \ge a) = P(S_n \ge a) + P(S_n \ge a+1) . \]
\item[(b)] Soit $ T_a = \min\{ k \ge 1 : S_k = a \} $ l'instant de premier passage en $ a \ge 1 $.
Démontrez le théorème du temps d'atteinte
\[ P(T_a = n) = \frac{a}{n}\, P(S_n = a) . \]
\item[(c)] Démontrez que la marche évite l'origine pendant $ 2n $ pas exactement aussi souvent qu'elle y
revient à l'instant $ 2n $ :
\[ P(S_1 \ne 0, S_2 \ne 0, \dots, S_{2n} \ne 0) = P(S_{2n} = 0) = u_{2n} , \]
et déduisez-en que le premier retour à l'origine a lieu à l'instant $ 2n $ avec probabilité
$ u_{2n}/(2n-1) $.
\item[(d)] Soit $ L_{2n} $ le nombre d'indices $ k \in \{1, \dots, 2n\} $ tels que
$ S_{k-1} \ge 0 $ et $ S_k \ge 0 $ --- le nombre de pas que la marche passe du côté positif.
Démontrez la loi de l'arcsinus discrète $ P(L_{2n} = 2k) = u_{2k}\, u_{2n-2k} $ pour
$ 0 \le k \le n $, et déduisez-en avec la formule de Stirling que
\[ P\!\left( \frac{L_{2n}}{2n} \le x \right) \longrightarrow \frac{2}{\pi}\arcsin\sqrt{x},
 \qquad 0 \le x \le 1 . \]
Expliquez ensuite ce que la forme de cette loi limite a de choquant.
\end{enumerate}

\end{problem}

\noindent Tout ceci est le problème du scrutin de PRB-36 devenu adulte : la réflexion
transforme des questions sur le chemin entier en questions sur son extrémité, le seul endroit où
l'on puisse lire un coefficient binomial. Feller a consacré à ces résultats un chapitre entier du
volume 1 et écrivait qu'ils contredisent l'intuition si violemment qu'aucune expérience n'y
prépare --- la densité de l'arcsinus en (d) est en forme de U, de sorte que dans une longue partie
équitable les issues les \textit{plus probables} sont celles où un joueur mène presque tout le temps,
un partage à peu près égal de l'avance étant la chose la plus rare qui soit. La partie (b) est
souvent appelée formule de Kemperman, et ses analogues continus sont les lois de l'arcsinus de Lévy
pour le mouvement brownien (PRB-42), démontrées par exactement la même réflexion.

\begin{refs}
\item Contexte : Principe de réflexion (processus de Wiener) --- Wikipédia (en anglais) (\url{https://en.wikipedia.org/wiki/Reflection_principle_(Wiener_process)}) (l'analogue continu)
\item Contexte : Lois de l'arcsinus (processus de Wiener) --- Wikipédia (en anglais) (\url{https://en.wikipedia.org/wiki/Arcsine_laws_(Wiener_process)})
\item Lecture : William Feller, \textit{An Introduction to Probability Theory and Its Applications}, vol. 1, ch. III
\item Lecture : Rick Durrett, \textit{Probability: Theory and Examples} (PDF gratuit, Duke (\url{https://sites.math.duke.edu/~rtd/PTE/PTE5_011119.pdf}))
\end{refs}

\bigskip


\section*{Master}

\begin{problem}[{Borel-Cantelli et la logique de l'\og{} infiniment souvent \fg{} --- Master}]
\textbf{PRB-15.} Pour des événements $ A_1, A_2, \dots $, on écrit
$ \{A_n \text{ i.s.}\} = \bigcap_m \bigcup_{n\ge m} A_n $ (\og{} $ A_n $ infiniment souvent \fg{}).
\begin{enumerate}
\item[(i)] Démontrez le premier lemme de Borel-Cantelli : $ \sum P(A_n) < \infty $ entraîne
$ P(A_n \text{ i.s.}) = 0 $.
\item[(ii)] Démontrez le second : si les $ A_n $ sont indépendants et $ \sum P(A_n) = \infty $, alors
$ P(A_n \text{ i.s.}) = 1 $ ; exhibez un contre-exemple montrant qu'on ne peut pas simplement
supprimer l'indépendance.
\item[(iii)] Application --- le singe à la machine à écrire : des lettres aléatoires i.i.d. uniformes
contiennent presque sûrement \og{} ABRACADABRA \fg{} une infinité de fois.
\item[(iv)] Application --- les plus longues séries : pour des lancers de pièce équilibrée, démontrez que la
plus longue série de \og{} face \fg{} $ L_n $ en $ n $ lancers vérifie $ L_n / \log_2 n \to 1 $
presque sûrement (majoration par le premier lemme, minoration par le second appliqué à des blocs
disjoints).
\end{enumerate}

\end{problem}

\noindent L'article de Borel de 1909 sur les nombres normaux et le raffinement de
Cantelli en 1917 ont donné aux probabilités leur mécanisme fondamental du zéro-un : les événements
asymptotiques ont tendance à avoir probabilité 0 ou 1, et ces deux lemmes sont ce qui permet de
trancher. La partie (iv) est la loi d'Erd\H{o}s-Rényi sous sa forme la plus simple --- la version
rigoureuse de l'analyse médico-légale des séries de PRB-05, et un premier aperçu de la façon dont
on démontre réellement les énoncés presque sûrs.

\begin{refs}
\item Contexte : Lemme de Borel-Cantelli --- Wikipédia (\url{https://fr.wikipedia.org/wiki/Lemme_de_Borel-Cantelli})
\item Lecture : Rick Durrett, \textit{Probability: Theory and Examples} (PDF gratuit, Duke (\url{https://sites.math.duke.edu/~rtd/PTE/PTE5_011119.pdf}))
\item Lecture : Geoffrey Grimmett \& David Stirzaker, \textit{Probability and Random Processes}
\end{refs}

\bigskip

\begin{problem}[{ABRACADABRA : martingales et arrêt optionnel --- Master}]
\textbf{PRB-16.} \begin{enumerate}
\item[(i)] Définissez une martingale et un temps d'arrêt ; énoncez le théorème d'arrêt optionnel, et
démontrez-le pour les temps d'arrêt bornés.
\item[(ii)] Un singe tape des majuscules i.i.d. uniformes. À chaque frappe, un nouveau joueur arrive, mise
1 sur le fait que la lettre suivante est un \og{} A \fg{}, et --- s'il est encore en vie --- laisse toute sa
fortune courir sur les lettres successives d'ABRACADABRA ; toutes les mises sont à cote équitable
(26 pour 1). En appliquant le théorème d'arrêt optionnel au profit total du casino, démontrez que
le nombre moyen de frappes avant la première apparition d'ABRACADABRA vaut
\[ 26^{11} + 26^4 + 26 , \]
et expliquez le rôle des auto-chevauchements du mot.
\item[(iii)] Déduisez-en la version pile ou face : le temps d'attente moyen vaut 10 pour FPF mais 8 pour
FFP --- démontrez les deux, et démontrez que parmi tous les motifs de trois lancers, ceux qui se
chevauchent eux-mêmes attendent toujours plus longtemps.
\item[(iv)] Vérifiez les hypothèses du théorème d'arrêt optionnel dans votre argument --- c'est là que
résident les vraies mathématiques.
\end{enumerate}

\end{problem}

\noindent L'argument de l'équipe de joueurs est celui de Shuo-Yen Robert Li (1980), qui
généralise des formules de Feller et les \og{} leading numbers \fg{} de Conway pour le jeu de Penney. C'est
la démonstration standard que l'arrêt optionnel des martingales n'est pas une abstraction mais un
instrument de calcul : une identité sur un temps d'attente où aucune martingale n'est visible tombe
en trois lignes une fois inventé le bon jeu équitable.

\begin{refs}
\item Contexte : Théorème d'arrêt de Doob --- Wikipédia (\url{https://fr.wikipedia.org/wiki/Th%C3%A9or%C3%A8me_d'arr%C3%AAt_de_Doob})
\item Article : S.-Y. R. Li, \og{} A Martingale Approach to the Study of Occurrence of Sequence Patterns in Repeated Experiments \fg{}, Ann. Probab. 8 (1980), Project Euclid (\url{https://projecteuclid.org/journals/annals-of-probability/volume-8/issue-6/A-Martingale-Approach-to-the-Study-of-Occurrence-of-Sequence/10.1214/aop/1176994578.full})
\item Lecture : David Williams, \textit{Probability with Martingales}
\end{refs}

\bigskip

\begin{problem}[{Percolation : une transition de phase que l'on peut démontrer --- Master}]
\textbf{PRB-17.} Dans la percolation par arêtes sur $ \mathbb{Z}^2 $,
chaque arête est \textit{ouverte} indépendamment avec probabilité $ p $. Posons
$ \theta(p) = P(\text{l'origine appartient à un amas ouvert infini}) $ et
$ p_c = \sup\{p : \theta(p) = 0\} $.
\begin{enumerate}
\item[(i)] Démontrez que $ \theta $ est croissante en $ p $ (construisez tous les processus de
percolation sur un même espace probabilisé).
\item[(ii)] Démontrez $ p_c \ge 1/3 $ : dénombrez les chemins auto-évitants issus de l'origine --- au plus
$ 4\cdot 3^{n-1} $ de longueur $ n $ --- et majorez la probabilité qu'un long chemin soit
ouvert.
\item[(iii)] Démontrez $ p_c \le 2/3 $ par l'argument de contour de Peierls : un amas fini est entouré
d'un circuit d'arêtes duales fermées ; dénombrez les circuits duaux et majorez leur probabilité
totale.
\item[(iv)] Concluez que $ 0 < p_c < 1 $ : une véritable transition de phase. Énoncez --- avec
attribution, sans démonstration --- la valeur exacte de $ p_c $ et qui l'a démontrée.
\end{enumerate}

\end{problem}

\noindent Broadbent et Hammersley inventèrent la percolation en 1957, pour modéliser la
façon dont un fluide --- ou le gaz de houille dans le filtre d'un masque --- pénètre un milieu poreux
aléatoire. Harris démontra $ \theta(1/2) = 0 $ en 1960 ; vingt ans plus tard, Kesten combla
l'écart avec $ p_c = 1/2 $ (1980), l'un des résultats célèbres des probabilités modernes. Les
deux arguments de dénombrement que vous donnez en (ii) et (iii) sont les outils originels du
domaine et restent la première chose que l'on essaie sur tout modèle nouveau ; la frontière à
laquelle ils mènent est PRB-20.

\begin{refs}
\item Contexte : Théorie de la percolation --- Wikipédia (\url{https://fr.wikipedia.org/wiki/Th%C3%A9orie_de_la_percolation})
\item Lecture : Geoffrey Grimmett, \textit{Percolation} (Springer)
\item Article : H. Duminil-Copin, \og{} Sixty years of percolation \fg{} (2017), arXiv:1712.04651 (\url{https://arxiv.org/abs/1712.04651})
\end{refs}

\bigskip

\begin{problem}[{Échangeabilité : le théorème de de Finetti --- Master}]
\textbf{PRB-18.} Une suite $ X_1, X_2, \dots $ est
\textit{échangeable} si sa loi est invariante par toute permutation d'un nombre fini d'indices.
\begin{enumerate}
\item[(i)] Montrez que toute suite i.i.d. est échangeable, et trouvez une suite échangeable qui n'est pas
i.i.d.
\item[(ii)] Démontrez le théorème de de Finetti pour les suites à valeurs dans $ \{0,1\} $ : une suite
échangeable infinie est un \textit{mélange} de suites i.i.d. de Bernoulli --- il existe une variable
aléatoire $ \Theta \in [0,1] $ telle que, conditionnellement à $ \Theta = \theta $, les
$ X_i $ soient i.i.d. de loi Bernoulli($ \theta $). (Une voie par les moments ou les
martingales : montrez que les fréquences empiriques convergent et conditionnez par la limite.)
\item[(iii)] Montrez que l'infinité est essentielle : exhibez une famille échangeable finie --- par exemple
un tirage sans remise dans une urne --- qui n'est en aucune façon un mélange de suites i.i.d.
\end{enumerate}

\end{problem}

\noindent Bruno de Finetti démontra le théorème en 1931 comme pierre angulaire de sa
théorie subjective des probabilités --- son mot d'ordre \og{} la probabilité n'existe pas \fg{} signifiait
que les chances objectives sont facultatives : pour qui croit à l'échangeabilité, le biais inconnu
$ \Theta $ est \textit{fabriqué} par le théorème, ce qui explique que la statistique bayésienne
traite les paramètres comme aléatoires. Hewitt et Savage (1955) l'ont étendu bien au-delà des
pièces de monnaie, et l'échangeabilité est devenue une théorie structurelle du hasard symétrique
(Aldous, Kingman) qui anime encore la recherche aujourd'hui.

\begin{refs}
\item Contexte : Théorème de de Finetti --- Wikipédia (en anglais) (\url{https://en.wikipedia.org/wiki/De_Finetti%27s_theorem})
\item Lecture : Rick Durrett, \textit{Probability: Theory and Examples} (PDF gratuit, Duke (\url{https://sites.math.duke.edu/~rtd/PTE/PTE5_011119.pdf}))
\item Lecture : Joseph K. Blitzstein \& Jessica Hwang, \textit{Introduction to Probability} (en libre accès (\url{https://probabilitybook.net}))
\end{refs}

\bigskip

\begin{problem}[{La formule de Wald --- Master, short}]
\textbf{PRB-29.} Soient $ X_1, X_2, \dots $ i.i.d. avec
$ E|X_1| < \infty $, soit $ N $ un temps d'arrêt pour la filtration qu'elles engendrent, avec
$ E[N] < \infty $, et posons $ S_N = X_1 + \cdots + X_N $. Démontrez la formule de Wald
$ E[S_N] = E[N]\, E[X_1] $, puis exhibez pour la marche aléatoire simple symétrique un temps
d'arrêt pour lequel l'identité est en défaut, en identifiant exactement quelle hypothèse a été
perdue.
\end{problem}

\noindent Abraham Wald démontra ceci alors qu'il dirigeait le programme d'analyse
séquentielle au Statistical Research Group de Columbia pendant la Seconde Guerre mondiale, où la
question était de savoir combien d'obus il faut inspecter avant d'accepter ou de rejeter un lot ;
c'est cette identité qui calcule la taille d'échantillon moyenne de son test séquentiel du rapport
de vraisemblance. C'est aussi le théorème d'arrêt optionnel de PRB-16 en habits ordinaires, et le
contre-exemple que vous construisez --- la marche arrêtée à son premier passage en $ +1 $, qui a
lieu avec probabilité un mais après un temps moyen infini --- est l'avertissement classique que
\og{} certain \fg{} et \og{} rapide \fg{} sont deux choses différentes.

\begin{refs}
\item Contexte : Formule de Wald --- Wikipédia (\url{https://fr.wikipedia.org/wiki/Formule_de_Wald})
\item Article : A. Wald, \og{} On cumulative sums of random variables \fg{}, Ann. Math. Statist. 15 (1944), 283--296
\item Lecture : David Williams, \textit{Probability with Martingales}
\end{refs}

\bigskip

\begin{problem}[{Une majoration exponentielle pour le pile ou face --- Master, short}]
\textbf{PRB-30.} Soit $ S_n $ le nombre de \og{} face \fg{} en $ n $ lancers d'une
pièce équilibrée. Démontrez la majoration de Chernoff-Hoeffding
\[ P\!\left( S_n \ge \tfrac{n}{2} + t \right) \le e^{-2t^2/n} \qquad (t > 0), \]
en montrant que les termes centrés vérifient $ E\big[e^{\theta (X_i - 1/2)}\big] \le
e^{\theta^2/8} $ pour tout réel $ \theta $, en appliquant l'inégalité de Markov à
$ e^{\theta S_n} $, puis en optimisant en $ \theta > 0 $.
\end{problem}

\noindent L'inégalité de Tchebychev (PRB-08) ne donne ici qu'une décroissance
polynomiale ; la méthode des moments exponentiels --- Bernstein dans les années 1920, Chernoff en
1952, Hoeffding en 1963 --- la remplace par une majoration qui décroît aussi vite que le théorème
central limite (PRB-14) le prévoit, et cela de manière non asymptotique. Cet écart entre
$ 1/k^2 $ et $ e^{-2t^2/n} $ est ce qui rend possibles les algorithmes randomisés, les codes
correcteurs d'erreurs et la théorie de l'apprentissage : c'est la raison pour laquelle quelques
centaines d'échantillons indépendants peuvent certifier une conclusion.

\begin{refs}
\item Contexte : Inégalité de Hoeffding --- Wikipédia (\url{https://fr.wikipedia.org/wiki/In%C3%A9galit%C3%A9_de_Hoeffding})
\item Contexte : Inégalité de Chernoff --- Wikipédia (\url{https://fr.wikipedia.org/wiki/In%C3%A9galit%C3%A9_de_Chernoff})
\item Article : W. Hoeffding, \og{} Probability inequalities for sums of bounded random variables \fg{}, J. Amer. Statist. Assoc. 58 (1963), 13--30
\end{refs}

\bigskip

\begin{problem}[{Extinction d'un processus de branchement --- Master}]
\textbf{PRB-31.} Dans un processus de Galton-Watson, $ Z_0 = 1 $ et chaque individu de la
génération $ n $ a, indépendamment, un nombre aléatoire d'enfants suivant la même loi de
reproduction $ (p_k)_{k \ge 0} $ ; notons $ f(s) = \sum_{k \ge 0} p_k s^k $ sa fonction
génératrice des probabilités et $ m = f'(1) = E[Z_1] $ le nombre moyen d'enfants. On suppose
$ p_1 < 1 $.
\begin{enumerate}
\item[(i)] Démontrez que la fonction génératrice de $ Z_n $ est la composée $ n $-ième
$ f_n = f \circ f \circ \cdots \circ f $, et déduisez-en $ E[Z_n] = m^n $.
\item[(ii)] Soit $ q = P(Z_n = 0 \text{ pour un certain } n) $ la probabilité d'extinction. Démontrez que
$ q $ est la plus petite racine de $ f(s) = s $ dans $ [0,1] $, et démontrez la dichotomie :
si $ m \le 1 $ alors $ q = 1 $, tandis que si $ m > 1 $ alors $ q < 1 $. (La convexité
de $ f $ sur $ [0,1] $ fait tout le travail.)
\item[(iii)] Démontrez que $ W_n = Z_n / m^n $ est une martingale positive, donc converge presque
sûrement ; énoncez ce qu'il faut de plus pour que la limite soit non dégénérée (la condition de
Kesten-Stigum) --- sans démonstration.
\item[(iv)] Calculez $ q $ explicitement lorsque la loi de reproduction est Poisson($ m $) puis
lorsqu'elle est géométrique, $ p_k = (1-a)a^k $. Diagnostiquez ensuite l'erreur historique :
Galton et Watson conclurent en 1874--75 que \textit{tout} patronyme finit par s'éteindre ; à l'aide de
(ii), dites exactement quelle racine ils ont prise et pourquoi la conclusion est fausse pour
$ m > 1 $.
\end{enumerate}

\end{problem}

\noindent Francis Galton demanda dans l'\textit{Educational Times} (1873) pourquoi tant de
familles distinguées avaient perdu leur patronyme, et Henry William Watson répondit avec des
fonctions génératrices --- bien posées, mal conclues, car il lut le point fixe $ s = 1 $ et en
déduisit l'extinction certaine. Irénée-Jules Bienaymé avait énoncé le critère correct dès 1845,
fait redécouvert seulement dans les années 1970 par Heyde et Seneta, et la lacune fut définitivement
comblée par Steffensen vers 1930. Le modèle a ensuite servi à calculer les réactions en chaîne de
neutrons à Los Alamos, et la condition $ m > 1 $ est exactement le $ R_0 > 1 $ des
épidémiologistes.

\begin{refs}
\item Contexte : Processus de Galton-Watson --- Wikipédia (\url{https://fr.wikipedia.org/wiki/Processus_de_Galton-Watson})
\item Contexte : Processus de branchement --- Wikipédia (\url{https://fr.wikipedia.org/wiki/Processus_de_branchement})
\item Article : H. W. Watson \& F. Galton, \og{} On the probability of the extinction of families \fg{}, J. Anthropological Institute of Great Britain and Ireland 4 (1875), 138--144
\item Lecture : Krishna B. Athreya \& Peter E. Ney, \textit{Branching Processes}
\end{refs}

\bigskip

\begin{problem}[{Conditionner, c'est projeter --- Master, short}]
\textbf{PRB-39.} Soient $ (\Omega, \mathcal{F}, P) $ un espace probabilisé,
$ \mathcal{G} \subseteq \mathcal{F} $ une sous-tribu, et
$ X \in L^2(\Omega, \mathcal{F}, P) $. Démontrez que $ L^2(\Omega, \mathcal{G}, P) $ est un
sous-espace fermé de $ L^2(\Omega, \mathcal{F}, P) $, que la projection orthogonale $ Y $ de
$ X $ sur ce sous-espace est caractérisée par
\[ E\big[ X \mathbf{1}_A \big] = E\big[ Y \mathbf{1}_A \big] \qquad \text{pour tout } A \in \mathcal{G}, \]
de sorte que $ Y = E[X \mid \mathcal{G}] $ est l'unique variable $ \mathcal{G} $-mesurable
minimisant $ E\big[(X - Y)^2\big] $ ; déduisez-en la propriété de la tour et la contraction
$ E\big[ E[X\mid\mathcal{G}]^2 \big] \le E[X^2] $ comme énoncés du théorème de Pythagore, et
expliquez pourquoi prolonger $ E[\cdot \mid \mathcal{G}] $ de $ L^2 $ à $ L^1 $ exige le
théorème de Radon-Nikodym plutôt que la géométrie.
\end{problem}

\noindent Kolmogorov définit l'espérance conditionnelle en 1933 comme une dérivée de
Radon-Nikodym précisément parce que le quotient élémentaire $ P(A \cap B)/P(B) $ s'effondre
lorsque $ P(B) = 0 $, et que conditionner par une observation continue revient toujours à
conditionner par des événements négligeables --- le paradoxe de Borel-Kolmogorov est ce qui arrive
quand on l'oublie. La lecture géométrique s'est révélée la plus féconde : les moindres carrés
(STA-04 sur la page Statistique (\url{applied-statistics.html})), la théorie de la prédiction
de Wiener dans les années 1940, le filtre de Kalman et tout le calcul des martingales sont la même
projection effectuée dans des espaces de Hilbert différents, et \og{} meilleur prédicteur \fg{} et
\og{} espérance conditionnelle \fg{} deviennent deux noms pour un seul objet.

\begin{refs}
\item Contexte : Espérance conditionnelle --- Wikipédia (\url{https://fr.wikipedia.org/wiki/Esp%C3%A9rance_conditionnelle})
\item Contexte : Théorème de projection sur un convexe fermé --- Wikipédia (\url{https://fr.wikipedia.org/wiki/Th%C3%A9or%C3%A8me_de_projection_sur_un_convexe_ferm%C3%A9})
\item Lecture : David Williams, \textit{Probability with Martingales}, ch. 9
\item Lecture : Rick Durrett, \textit{Probability: Theory and Examples}, ch. 4 (PDF gratuit, Duke (\url{https://sites.math.duke.edu/~rtd/PTE/PTE5_011119.pdf}))
\end{refs}

\bigskip

\begin{problem}[{Fonctions caractéristiques et théorème de convergence de Lévy --- Master}]
\textbf{PRB-40.} Pour une variable aléatoire réelle $ X $, on définit sa fonction
caractéristique $ \varphi_X(t) = E\big[e^{itX}\big] $, $ t \in \mathbb{R} $.
\begin{enumerate}
\item[(a)] Démontrez que $ \varphi_X $ est uniformément continue, vérifie $ \varphi_X(0) = 1 $ et
$ |\varphi_X(t)| \le 1 $, que $ \varphi_{X+Y} = \varphi_X \varphi_Y $ pour $ X, Y $
indépendantes, et que si $ E|X|^k $ est fini alors $ \varphi_X $ est $ k $ fois dérivable
avec $ \varphi_X^{(k)}(0) = i^k E[X^k] $.
\item[(b)] Démontrez la formule d'inversion de Lévy : pour $ a < b $,
\[ \lim_{T \to \infty} \frac{1}{2\pi} \int_{-T}^{T}
 \frac{e^{-ita} - e^{-itb}}{it}\, \varphi_X(t)\, dt
 = P(a < X < b) + \tfrac{1}{2}\big( P(X = a) + P(X = b) \big) , \]
et déduisez-en le théorème d'unicité : la fonction caractéristique détermine la loi.
\item[(c)] Énoncez le théorème de convergence de Lévy --- si $ \varphi_{X_n}(t) \to \varphi(t) $ pour tout
$ t $ et si $ \varphi $ est continue en $ 0 $, alors $ \varphi $ est la fonction
caractéristique d'une certaine variable $ X $ et $ X_n $ converge en loi vers $ X $ --- et
démontrez le sens facile. Démontrez ensuite le théorème central limite pour des variables i.i.d.
d'espérance $ \mu $ et de variance finie $ \sigma^2 $ en développant $ \log \varphi $ à
l'ordre deux, et dites exactement où la continuité de la limite en $ 0 $ intervient.
\item[(d)] Montrez pourquoi la méthode des moments ne peut pas remplacer tout cela. Démontrez que pour la
densité log-normale standard $ f $ sur $ (0,\infty) $, les densités perturbées
\[ f_c(x) = f(x)\big( 1 + c \sin(2\pi \log x) \big), \qquad |c| \le 1 , \]
sont de véritables densités de probabilité ayant exactement les mêmes moments de tous ordres que
$ f $ --- une loi n'est donc pas nécessairement déterminée par ses moments, alors que d'après (b)
elle l'est toujours par sa fonction caractéristique.
\end{enumerate}

\end{problem}

\noindent Laplace et Cauchy utilisaient déjà les transformées de Fourier de lois, mais
c'est Liapounov (1901) qui en fit l'instrument standard des théorèmes limites, et Lévy qui bâtit la
théorie sous la forme employée ici ; le théorème de convergence est ce qui transforme un calcul sur
des fonctions de $ t $ en un énoncé de convergence de lois, et toute démonstration moderne du
théorème central limite passe par lui. La partie (d) est la découverte de Stieltjes en 1894 sous la
forme donnée par C. C. Heyde en 1963 : la loi log-normale, l'une des plus utilisées en statistique,
partage ses moments avec toute une famille d'autres lois. C'est la raison précise pour laquelle la
méthode des moments (STA-22 sur la page Statistique (\url{applied-statistics.html})) a besoin
d'une condition auxiliaire comme celle de Carleman, et pourquoi les fonctions caractéristiques ont
gagné.

\begin{refs}
\item Contexte : Fonction caractéristique (probabilités) --- Wikipédia (\url{https://fr.wikipedia.org/wiki/Fonction_caract%C3%A9ristique_(probabilit%C3%A9s)})
\item Contexte : Théorème de convergence de Lévy --- Wikipédia (\url{https://fr.wikipedia.org/wiki/Th%C3%A9or%C3%A8me_de_convergence_de_L%C3%A9vy})
\item Article : C. C. Heyde, \og{} On a property of the lognormal distribution \fg{}, J. Roy. Statist. Soc. B 25 (1963), 392--393
\item Lecture : Rick Durrett, \textit{Probability: Theory and Examples}, ch. 3 (PDF gratuit, Duke (\url{https://sites.math.duke.edu/~rtd/PTE/PTE5_011119.pdf}))
\end{refs}

\bigskip

\begin{problem}[{Le théorème de Cramér : le prix d'une grande déviation --- Master}]
\textbf{PRB-41.} Soient $ X_1, X_2, \dots $ i.i.d. telles que
$ \Lambda(\theta) = \log E\big[e^{\theta X_1}\big] $ soit finie pour tout réel $ \theta $,
d'espérance $ \mu = E[X_1] $, et posons $ S_n = X_1 + \cdots + X_n $. On définit la fonction de
taux comme la transformée de Legendre
\[ I(x) = \sup_{\theta \in \mathbb{R}} \big( \theta x - \Lambda(\theta) \big) . \]
\begin{enumerate}
\item[(a)] Démontrez que $ \Lambda $ est convexe et dérivable, que $ I $ est convexe et positive, et
que $ I(\mu) = 0 $.
\item[(b)] Démontrez la majoration non asymptotique $ P(S_n \ge nx) \le e^{-n I(x)} $ pour tout
$ x \ge \mu $ et tout $ n $, en optimisant l'inégalité de Markov appliquée à
$ e^{\theta S_n} $ (c'est PRB-30 sous forme générale).
\item[(c)] Démontrez la minoration correspondante par bascule exponentielle : pour $ x $ à l'intérieur
de l'image de $ \Lambda' $, soit $ \theta^{*} $ la solution de
$ \Lambda'(\theta^{*}) = x $, définissez la loi basculée
$ dP_{\theta^{*}} = e^{\theta^{*} y - \Lambda(\theta^{*})} dP $, et montrez que sous cette loi la
somme a pour espérance $ nx $, de sorte qu'une estimation de type théorème central limite donne
\[ \lim_{n \to \infty} \frac{1}{n} \log P(S_n \ge nx) = -I(x) . \]
\item[(d)] Calculez $ I $ explicitement pour une pièce équilibrée ($ X_i \in \{0,1\} $) et pour la loi
exponentielle de taux 1, et vérifiez que $ I(x) \approx (x-\mu)^2/(2\sigma^2) $ pour $ x $
proche de $ \mu $. Montrez ensuite sur un exemple qu'extrapoler l'approximation gaussienne de
PRB-14 dans le régime des grandes déviations donne le mauvais exposant --- le théorème central limite
décrit une fenêtre de largeur $ \sqrt{n} $ et ne dit rien des événements de probabilité
$ e^{-cn} $.
\end{enumerate}

\end{problem}

\noindent Harald Cramér démontra ceci en 1938 alors qu'il travaillait sur la théorie
collective du risque de l'assurance suédoise --- la question y est la probabilité que les réserves
d'une compagnie soient ruinées par une série improbable de sinistres, et c'est précisément une
probabilité que la loi des grands nombres déclare petite sans dire à quel point. Chernoff
redécouvrit la majoration en 1952 pour des problèmes de test ; la formulation du principe de
grandes déviations par Varadhan en 1966 (prix Abel 2007) a transformé le motif en théorie générale,
étendue par Sanov aux mesures empiriques et par Freidlin et Wentzell aux petites perturbations
aléatoires de systèmes dynamiques. La morale de la partie (d) est à retenir : une approximation
gaussienne est un énoncé sur le centre d'une loi, et la queue est gouvernée par une tout autre
fonction.

\begin{refs}
\item Contexte : Théorème de Cramér --- Wikipédia (\url{https://fr.wikipedia.org/wiki/Th%C3%A9or%C3%A8me_de_Cram%C3%A9r})
\item Contexte : Principe de grandes déviations --- Wikipédia (\url{https://fr.wikipedia.org/wiki/Principe_de_grandes_d%C3%A9viations})
\item Article : H. Cramér, \og{} Sur un nouveau théorème-limite de la théorie des probabilités \fg{}, Actualités Scientifiques et Industrielles 736 (1938)
\item Lecture : Amir Dembo \& Ofer Zeitouni, \textit{Large Deviations Techniques and Applications}, ch. 2
\end{refs}

\bigskip

\begin{problem}[{Construire le mouvement brownien, et démontrer qu'il n'est nulle part régulier --- Master}]
\textbf{PRB-42.} Un mouvement brownien standard sur $ [0,1] $ est un processus
$ (B_t)_{0 \le t \le 1} $ vérifiant $ B_0 = 0 $, à accroissements indépendants, avec
$ B_t - B_s \sim N(0, t-s) $ pour $ s < t $, et à trajectoires continues. Rien ne garantit
d'avance qu'un tel objet existe.
\begin{enumerate}
\item[(a)] Construisez-le. Soit $ (h_k) $ la base orthonormée de Haar de $ L^2[0,1] $, soit
$ \Delta_k(t) = \int_0^t h_k(s)\,ds $ les fonctions de Schauder, et soient $ \xi_k $ des
variables gaussiennes standard i.i.d. Démontrez que
\[ B_t = \sum_{k \ge 0} \xi_k \, \Delta_k(t) \]
converge uniformément sur $ [0,1] $ avec probabilité un --- majorez $ \max |\xi_k| $ sur chaque
bloc dyadique et utilisez Borel-Cantelli (PRB-15) --- et vérifiez que la limite possède les trois
propriétés ci-dessus.
\item[(b)] Démontrez que les trajectoires sont presque sûrement höldériennes de tout exposant
$ \alpha < 1/2 $, et démontrez qu'elles ne sont presque sûrement \textit{pas} höldériennes
d'exposant $ 1/2 $.
\item[(c)] Démontrez qu'avec probabilité un la trajectoire $ t \mapsto B_t $ n'est dérivable en aucun
point. (Si $ B $ avait une dérivée en un certain $ s $, alors pour $ n $ grand plusieurs
accroissements consécutifs de taille $ 2^{-n} $ près de $ s $ seraient tous anormalement
petits ; majorez la probabilité de cet événement et sommez sur la grille dyadique.)
\item[(d)] Démontrez que le long des partitions dyadiques la variation quadratique
$ \sum_{j} \big( B_{(j+1)2^{-n}} - B_{j 2^{-n}} \big)^2 $ converge vers $ t $ dans
$ L^2 $, et déduisez-en que les trajectoires sont à variation totale infinie sur tout intervalle.
Expliquez en un paragraphe pourquoi c'est la raison précise pour laquelle $ \int f\, dB $ ne peut
pas être une intégrale de Stieltjes, et pourquoi Itô a dû en inventer une autre.
\end{enumerate}

\end{problem}

\noindent Robert Brown observa des grains de pollen frémir en 1827 ; Bachelier modélisa
la Bourse de Paris avec le même processus en 1900, et Einstein s'en servit en 1905 pour plaider la
réalité des atomes --- mais personne n'avait produit d'objet mathématique possédant ces propriétés
avant que Norbert Wiener ne le fasse en 1923, d'où le nom de processus de Wiener. La construction
de (a) est celle de Lévy (1948) sous la forme en fonctions de Haar de Ciesielski (1961), et c'est
la démonstration d'existence la plus propre que l'on connaisse : une série aléatoire explicite de
type Fourier. La partie (c) est le théorème de Paley, Wiener et Zygmund (1933), qui a choqué les
analystes --- pendant un siècle les mathématiciens avaient tenu la fonction nulle part dérivable de
Weierstrass pour un monstre, et voici que presque toute trajectoire du processus le plus naturel
des probabilités en est une. La partie (d) est le point de départ du calcul stochastique : la
variation quadratique non nulle est exactement le terme supplémentaire de la formule d'Itô.

\begin{refs}
\item Contexte : Processus de Wiener --- Wikipédia (\url{https://fr.wikipedia.org/wiki/Processus_de_Wiener})
\item Contexte : Mouvement brownien --- Wikipédia (\url{https://fr.wikipedia.org/wiki/Mouvement_brownien})
\item Article : R. E. A. C. Paley, N. Wiener \& A. Zygmund, \og{} Notes on random functions \fg{}, Math. Z. 37 (1933), 647--668
\item Lecture : Peter Mörters \& Yuval Peres, \textit{Brownian Motion} (Cambridge, 2010), ch. 1
\end{refs}

\bigskip


\section*{Recherche}

\begin{problem}[{Chemins auto-évitants et constante de connectivité --- Recherche}]
\textbf{PRB-19.} Soit $ c_n $ le nombre de chemins auto-évitants de
$ n $ pas issus de l'origine dans un réseau.
\begin{enumerate}
\item[(i)] Démontrez --- cette partie est entièrement à votre portée --- que $ c_{m+n} \le c_m c_n $, et
déduisez-en, via le lemme de sous-additivité de Fekete (énoncez-le et démontrez-le), que la
\textit{constante de connectivité} $ \mu = \lim_n c_n^{1/n} $ existe.
\item[(ii)] Démontrez les majorations et minorations grossières $ 2 \le \mu \le 3 $ pour
$ \mathbb{Z}^2 $.
\item[(iii)] \textbf{Résolu (2012) :} sur le réseau en nid d'abeilles,
$ \mu = \sqrt{2 + \sqrt{2}} $ --- prédit par le physicien Nienhuis en 1982, démontré par
Duminil-Copin et Smirnov à l'aide d'une observable discrètement holomorphe. Comprenez l'énoncé et
l'architecture de la démonstration.
\item[(iv)] \textbf{Ouvert :} la valeur exacte de $ \mu $ pour $ \mathbb{Z}^2 $ (numériquement
$\approx$ 2,63816) est inconnue ; le comportement asymptotique conjecturé
$ c_n \sim A\,\mu^n n^{11/32} $ dans le plan et la conjecture selon laquelle la limite d'échelle
est l'évolution de Schramm-Loewner SLE8/3 ne sont pas démontrés ; en dimension 3,
essentiellement rien n'est rigoureux sur les exposants critiques.
\end{enumerate}

\end{problem}

\noindent Le chemin auto-évitant est le modèle standard d'une longue chaîne polymère,
introduit par le chimiste Flory dans les années 1940 ; il est notoirement difficile parce que la
contrainte d'auto-évitement détruit l'indépendance. Qu'un seul réseau --- celui en nid d'abeilles ---
ait livré une constante exacte, grâce à des idées importées de la théorie conforme des champs, est
l'un des résultats célèbres des probabilités du XXIe siècle ; toutes les questions
voisines restent ouvertes, énoncées honnêtement ci-dessus.

\begin{refs}
\item Contexte : Chemin auto-évitant --- Wikipédia (\url{https://fr.wikipedia.org/wiki/Chemin_auto-%C3%A9vitant})
\item Contexte : Constante de connectivité --- Wikipédia (\url{https://fr.wikipedia.org/wiki/Constante_de_connectivit%C3%A9})
\item Article : H. Duminil-Copin \& S. Smirnov, \og{} The connective constant of the honeycomb lattice equals $\surd$(2+$\surd$2) \fg{}, Ann. of Math. 175 (2012), arXiv:1007.0575 (\url{https://arxiv.org/abs/1007.0575})
\item Lecture : Neal Madras \& Gordon Slade, \textit{The Self-Avoiding Walk}
\end{refs}

\bigskip

\begin{problem}[{Percolation critique en dimension trois --- Recherche}]
\textbf{PRB-20.} Poursuivez à partir de PRB-17.
\begin{enumerate}
\item[(i)] Dans le plan, le tableau critique est un théorème : pour la percolation par sites sur le
réseau triangulaire, Smirnov a démontré l'invariance conforme des probabilités de traversée
(formule de Cardy, 2001) et, avec Werner, en a déduit les exposants critiques --- par exemple
$ \theta(p) = (p - p_c)^{5/36 + o(1)} $ lorsque $ p \downarrow p_c $. Comprenez ces énoncés et
la façon dont SLE6 intervient.
\item[(ii)] En grande dimension ($ d \ge 11 $, Fitzner-van der Hofstad 2017, après Hara-Slade pour
$ d \ge 19 $), le développement en lacets donne les exposants de champ moyen.
\item[(iii)] \textbf{Ouvert :} pour $ 3 \le d \le 10 $ --- y compris la dimension physique 3 --- il n'est pas
démontré que les exposants critiques existent, et encore moins quelles sont leurs valeurs ; même la
continuité de la transition de phase (à savoir $ \theta(p_c) = 0 $) est un célèbre problème
ouvert dans ces dimensions. Tâche accessible : démontrez que $ \theta(p_c) = 0 $
\textit{découlerait}, pour $ d = 3 $, d'hypothèses d'unicité de la limite d'échelle que vous
énoncerez précisément ; lisez ensuite ce qui est réellement connu.
\end{enumerate}

\end{problem}

\noindent Les physiciens disposent de valeurs numériques précises pour les exposants en
dimension 3 et ne doutent pas de leur existence ; les mathématiques ne savent pas encore démontrer
cette existence. La solution en dimension deux (Smirnov, médaille Fields 2010) et le développement
en lacets en grande dimension encadrent le mystère sans y toucher --- la dimension trois est la
frontière où les probabilités, la physique statistique et la géométrie conforme s'arrêtent
aujourd'hui. État en 2026 : ouvert. Voir aussi
Problèmes ouverts (\url{open-problems.html}).

\begin{refs}
\item Contexte : Exposants critiques de la percolation --- Wikipédia (en anglais) (\url{https://en.wikipedia.org/wiki/Percolation_critical_exponents})
\item Article : S. Smirnov, \og{} Critical percolation in the plane: conformal invariance, Cardy's formula, scaling limits \fg{}, C. R. Acad. Sci. Paris 333 (2001)
\item Article : H. Duminil-Copin, \og{} Sixty years of percolation \fg{} (2017), arXiv:1712.04651 (\url{https://arxiv.org/abs/1712.04651})
\item Lecture : Geoffrey Grimmett, \textit{Percolation} (Springer)
\end{refs}

\bigskip

\begin{problem}[{L'universalité KPZ --- Recherche}]
\textbf{PRB-21.} De nombreux modèles de croissance unidimensionnels ---
croissance en coin, percolation de dernier passage, processus d'exclusion asymétrique, taches de
café s'étalant sur du papier --- sont conjecturés partager une même limite d'échelle : des
fluctuations de hauteur d'ordre $ t^{1/3} $, des corrélations sur une distance $ t^{2/3} $, et
des lois limites de Tracy-Widom issues de la théorie des matrices aléatoires.
\begin{enumerate}
\item[(i)] Parties accessibles : dérivez heuristiquement le changement d'échelle 1:2:3 à partir de
l'équation KPZ $ \partial_t h = \tfrac12 \partial_x^2 h + \tfrac12 (\partial_x h)^2 +
\xi $, et comprenez pourquoi le TCL gaussien (PRB-14) \textit{ne peut pas} valoir ici.
\item[(ii)] \textbf{Résolu pour les modèles exactement solubles :} le processus de Markov limite --- le point
fixe KPZ --- a été construit par Matetski, Quastel et Remenik (Acta Math. 227, 2021) à partir du
TASEP, et le paysage dirigé, plus riche, par Dauvergne, Ortmann et Virág (2018--22) à partir de la
percolation de dernier passage brownienne ; on sait désormais que l'équation KPZ elle-même converge
vers le point fixe. Comprenez ces énoncés.
\item[(iii)] \textbf{Ouvert :} l'universalité --- le fait que les modèles génériques de la classe (lois de poids
générales, règles microscopiques quelconques) convergent vers la même limite --- reste indémontrée
hors de l'îlot intégrable.
\end{enumerate}

\end{problem}

\noindent Kardar, Parisi et Zhang écrivirent leur équation en 1986 ; le sujet explosa
quand Baik, Deift et Johansson (1999) trouvèrent des fluctuations de Tracy-Widom dans les plus
longues sous-suites croissantes. Les constructions des années 2020 ont transformé une prophétie de
physiciens en objets mathématiques, mais la \textit{conjecture} d'universalité --- l'affirmation que le
point fixe mérite son nom --- est ouverte, et largement tenue pour l'un des problèmes centraux des
probabilités modernes. État en 2026 : ouvert au-delà des modèles exactement solubles.

\begin{refs}
\item Contexte : Équation de Kardar-Parisi-Zhang --- Wikipédia (en anglais) (\url{https://en.wikipedia.org/wiki/Kardar%E2%80%93Parisi%E2%80%93Zhang_equation})
\item Contexte : Loi de Tracy-Widom --- Wikipédia (en anglais) (\url{https://en.wikipedia.org/wiki/Tracy%E2%80%93Widom_distribution})
\item Article : K. Matetski, J. Quastel \& D. Remenik, \og{} The KPZ fixed point \fg{}, Acta Math. 227 (2021), Project Euclid (\url{https://projecteuclid.org/journals/acta-mathematica/volume-227/issue-1/The-KPZ-fixed-point/10.4310/ACTA.2021.v227.n1.a3.full})
\item Article : D. Dauvergne, J. Ortmann \& B. Virág, \og{} The directed landscape \fg{} (2018), arXiv:1812.00309 (\url{https://arxiv.org/abs/1812.00309})
\end{refs}

\bigskip

\begin{problem}[{Ensembles de boucles conformes : les boucles des modèles critiques --- Recherche}]
\textbf{PRB-22.} L'ensemble de boucles conforme CLE$\kappa$
($ 8/3 < \kappa < 8 $) est une collection aléatoire de boucles sans croisement dans un
domaine plan, dont la loi est invariante par les applications conformes ; c'est la limite d'échelle
conjecturée de toute la famille des interfaces des modèles critiques bidimensionnels sur
réseau.
\begin{enumerate}
\item[(i)] Parties accessibles : précisez ce que signifie \og{} limite d'échelle conformément invariante \fg{}
pour un modèle sur réseau, et vérifiez l'invariance conforme dans le seul cas où elle est
élémentaire --- la trajectoire du mouvement brownien plan (Lévy).
\item[(ii)] \textbf{Cas résolus, avec attribution :} les interfaces de la percolation critique par sites
convergent vers CLE6 (Smirnov 2001 ; Camia-Newman 2006), et les interfaces d'Ising
critique vers CLE3 (Benoist-Hongler, Ann. Probab. 47, 2019) ; Sheffield et Werner
(Ann. of Math. 176, 2012) ont caractérisé CLE axiomatiquement. Comprenez ces énoncés.
\item[(iii)] \textbf{Ouvert :} la convergence vers CLE$\kappa$ pour la grande majorité des modèles
conjecturés --- chemin auto-évitant vers SLE8/3 ($\kappa$ = 8/3), modèles de boucles O(n) pour
n général, variantes de l'arbre couvrant uniforme au-delà des cas démontrés --- reste indémontrée ;
chacun de ces cas est un problème de frontière identifié.
\end{enumerate}

\end{problem}

\noindent L'invention de SLE par Schramm (1999) a transformé la théorie conforme des
champs des physiciens en théorèmes de probabilités et réorganisé tout le domaine des phénomènes
critiques bidimensionnels ; CLE en est le complément à plusieurs boucles. Le tableau du domaine est
net : l'universalité est crue pour tous les modèles et démontrée seulement là où l'on a trouvé une
miraculeuse holomorphie discrète --- exactement la situation de PRB-19 et PRB-20, et la raison pour
laquelle les phénomènes critiques plans restent une mine d'or en activité. État en 2026 : ouvert
hors des cas listés.

\begin{refs}
\item Contexte : Ensemble de boucles conforme --- Wikipédia (en anglais) (\url{https://en.wikipedia.org/wiki/Conformal_loop_ensemble})
\item Contexte : Évolution de Schramm-Loewner --- Wikipédia (en anglais) (\url{https://en.wikipedia.org/wiki/Schramm%E2%80%93Loewner_evolution})
\item Article : S. Sheffield \& W. Werner, \og{} Conformal loop ensembles: the Markovian characterization and the loop-soup construction \fg{}, Ann. of Math. 176 (2012)
\item Article : S. Benoist \& C. Hongler, \og{} The scaling limit of critical Ising interfaces is CLE(3) \fg{}, Ann. Probab. 47 (2019), arXiv:1604.06975 (\url{https://arxiv.org/abs/1604.06975})
\end{refs}

\bigskip

\begin{problem}[{Énoncez précisément ce que la formule de Parisi a réglé --- Recherche, short}]
\textbf{PRB-32.} Le verre de spin de Sherrington-Kirkpatrick attribue à chaque
configuration $ \sigma \in \{-1,+1\}^N $ l'énergie
$ H_N(\sigma) = N^{-1/2} \sum_{i < j} g_{ij}\, \sigma_i \sigma_j $ avec des couplages
gaussiens standard i.i.d. $ g_{ij} $, et l'ansatz de brisure de symétrie des répliques de Parisi
(1979--80) prédisait la limite de $ N^{-1} E \log \sum_{\sigma} e^{\beta H_N(\sigma)} $. Rédigez un
compte rendu précis de l'état des lieux : énoncez la formule variationnelle de Parisi, énoncez
exactement ce que Guerra (2003), Talagrand (2006) et le théorème d'ultramétricité de Panchenko
(2013) ont chacun établi, et énoncez honnêtement que le tableau correspondant pour le modèle
d'Edwards-Anderson en dimension finie reste ouvert.
\end{problem}

\noindent La solution de Parisi fut obtenue par un calcul faisant intervenir
\og{} $ n \to 0 $ répliques \fg{} qui n'avait absolument aucun sens mathématique, et pendant deux
décennies on ignora si c'était même la bonne réponse. L'interpolation de Guerra fournit une
inégalité, Talagrand l'autre ; la part de Parisi dans le prix Nobel de physique 2021 citait ce
cercle d'idées. Ce que l'on ne sait \textit{pas}, c'est si un tableau semblable vaut pour des verres
de spin réalistes à courte portée en trois dimensions --- la controverse entre le modèle des
gouttelettes et la brisure de symétrie des répliques n'est pas tranchée. État en 2026 : cas de
champ moyen démontré, cas en dimension finie ouvert.

\begin{refs}
\item Contexte : Modèle de Sherrington-Kirkpatrick --- Wikipédia (en anglais) (\url{https://en.wikipedia.org/wiki/Sherrington%E2%80%93Kirkpatrick_model})
\item Contexte : Verre de spin --- Wikipédia (\url{https://fr.wikipedia.org/wiki/Verre_de_spin})
\item Article : M. Talagrand, \og{} The Parisi formula \fg{}, Ann. of Math. 163 (2006), 221--263
\item Article : D. Panchenko, \og{} The Parisi ultrametricity conjecture \fg{}, Ann. of Math. 177 (2013), 383--393
\end{refs}

\bigskip

\begin{problem}[{Où le 3-SAT aléatoire devient-il insatisfiable ? --- Recherche, short}]
\textbf{PRB-33.} Une formule $ k $-SAT aléatoire tire $ m = \alpha n $
clauses uniformément et indépendamment parmi les $ \binom{n}{k} 2^k $ $ k $-clauses possibles
sur $ n $ variables booléennes, et la conjecture du seuil de satisfaisabilité affirme que pour
chaque $ k $ il existe un $ \alpha_k $ tel que la formule soit satisfiable avec une probabilité
tendant vers 1 en dessous de $ \alpha_k $ et vers 0 au-dessus. Rédigez un compte rendu précis de
l'état des lieux : Friedgut (1999) a démontré l'existence d'une \textit{suite} de seuils nets sans la
localiser, Ding, Sly et Sun (Ann. of Math. 196, 2022) ont déterminé $ \alpha_k $ pour tout $ k $
assez grand, et le cas $ k = 3 $ --- où la méthode de la cavité des physiciens prédit
$ \alpha_3 \approx 4{,}267 $ --- reste ouvert.
\end{problem}

\noindent C'est le point de rencontre le plus net entre probabilités, physique statistique
et informatique théorique : la formule conjecturée pour $ \alpha_k $ est issue de la méthode de
la cavité de Mézard, Parisi et Zecchina, la même machinerie de répliques qu'en PRB-32, et la
démonstration pour $ k $ grand a dû transformer ce calcul non rigoureux en théorème.
Empiriquement, la transition est spectaculaire --- au voisinage de $ \alpha_3 \approx 4{,}267 $ les
formules aléatoires deviennent les plus difficiles pour les solveurs SAT, ce qui rend le seuil
intéressant en pratique autant qu'en théorie. État en 2026 : ouvert pour $ k = 3 $.

\begin{refs}
\item Contexte : Problème SAT --- Wikipédia (\url{https://fr.wikipedia.org/wiki/Probl%C3%A8me_SAT})
\item Article : E. Friedgut, \og{} Sharp thresholds of graph properties, and the $ k $-SAT problem \fg{}, J. Amer. Math. Soc. 12 (1999), 1017--1054
\item Article : J. Ding, A. Sly \& N. Sun, \og{} Proof of the satisfiability conjecture for large $ k $ \fg{}, Ann. of Math. 196 (2022), arXiv:1411.0650 (\url{https://arxiv.org/abs/1411.0650})
\end{refs}

\bigskip

\begin{problem}[{Temps de mélange et phénomène de coupure --- Recherche}]
\textbf{PRB-34.} Pour une chaîne de Markov irréductible apériodique sur un espace d'états fini,
de matrice de transition $ P $ et de loi stationnaire $ \pi $, posons
\[ d(t) = \max_{x} \; \big\lVert P^t(x, \cdot) - \pi \big\rVert_{\mathrm{TV}} ,
\qquad t_{\mathrm{mix}}(\varepsilon) = \min\{ t : d(t) \le \varepsilon \} . \]
Une \textit{famille} de chaînes indexée par $ n $ présente une \textit{coupure} (cutoff) si
$ t_{\mathrm{mix}}^{(n)}(\varepsilon) / t_{\mathrm{mix}}^{(n)}(1-\varepsilon) \to 1 $ pour tout
$ \varepsilon \in (0, 1/2) $ --- la distance à l'équilibre reste proche de 1 puis s'effondre vers 0
dans une fenêtre d'ordre inférieur au temps de mélange lui-même.
\begin{enumerate}
\item[(i)] Partie accessible : démontrez que $ d $ est décroissante et que
$ \bar{d}(t) = \max_{x,y} \lVert P^t(x,\cdot) - P^t(y,\cdot) \rVert_{\mathrm{TV}} $ est
sous-multiplicative, $ \bar{d}(s+t) \le \bar{d}(s)\bar{d}(t) $, et déduisez-en que la distance à
la stationnarité décroît exponentiellement une fois $ t_{\mathrm{mix}} $ dépassé.
\item[(ii)] Démontrez une coupure de bout en bout pour la marche aléatoire paresseuse sur l'hypercube
$ \{0,1\}^n $ --- à chaque pas, on choisit une coordonnée uniformément et on la remplace par un
nouveau tirage à pile ou face équilibré : montrez que
$ t_{\mathrm{mix}} = \tfrac12 n \log n\,(1+o(1)) $, la minoration venant du dénombrement à la
collectionneur de vignettes des coordonnées jamais touchées (PRB-07) et la majoration d'un
couplage.
\item[(iii)] Comprenez le cas célèbre : Bayer et Diaconis (1992) ont démontré que le battage américain
(riffle shuffle) d'un jeu de $ n $ cartes présente une coupure à $ \tfrac32 \log_2 n $
battages, source des fameux \og{} sept battages \fg{} pour $ n = 52 $. Énoncez précisément ce que ce
théorème affirme et ce qu'il n'affirme pas.
\item[(iv)] \textbf{Ouvert :} Aldous et Diaconis ont demandé dans les années 1980 un critère général de
coupure. La \textit{condition de produit} de Peres --- le produit de $ t_{\mathrm{mix}} $ par le trou
spectral tend vers l'infini --- est nécessaire, mais le contre-exemple d'Aldous montre qu'elle n'est
pas suffisante en général ; elle l'\textit{est} pour les chaînes de naissance et de mort
(Ding-Lubetzky-Peres, 2010) et pour les chaînes à courbure positive (Salez, 2021). Savoir si elle
suffit pour des familles symétriques naturelles --- par exemple toutes les chaînes sommet-transitives
--- est ouvert. Énoncez la conjecture précisément et lisez ce qui est connu.
\end{enumerate}

\end{problem}

\noindent La coupure fut découverte par Aldous et Diaconis vers 1981 en analysant des
battages de cartes : pendant longtemps rien ne semble se passer, puis le jeu est soudain aléatoire.
Le phénomène est aujourd'hui connu pour des dizaines de chaînes et reste invisible aux majorations
spectrales grossières, de sorte que chaque démonstration a besoin d'une statistique discriminante
pour la minoration et d'un couplage ou d'un temps stationnaire fort pour la majoration. Qu'il
n'existe toujours pas de critère général --- après quarante ans, avec la condition de produit si
proche du but --- est l'un des problèmes ouverts bien connus du domaine. État en 2026 : ouvert en
général, démontré pour les familles particulières listées ci-dessus.

\begin{refs}
\item Contexte : Temps de mélange d'une chaîne de Markov --- Wikipédia (en anglais) (\url{https://en.wikipedia.org/wiki/Markov_chain_mixing_time})
\item Lecture : David A. Levin \& Yuval Peres, \textit{Markov Chains and Mixing Times}, 2e éd. (AMS)
\item Article : D. Bayer \& P. Diaconis, \og{} Trailing the dovetail shuffle to its lair \fg{}, Ann. Appl. Probab. 2 (1992), 294--313
\item Articles : J. Ding, E. Lubetzky \& Y. Peres, \og{} Total variation cutoff in birth-and-death chains \fg{}, Probab. Theory Related Fields 146 (2010), 61--85 ; J. Salez, \og{} Cutoff for non-negatively curved Markov chains \fg{}, J. Eur. Math. Soc., arXiv:2102.05597 (\url{https://arxiv.org/abs/2102.05597})
\end{refs}

\bigskip

\begin{problem}[{L'inégalité de corrélation gaussienne --- Recherche, short}]
\textbf{PRB-43.} L'inégalité de corrélation gaussienne affirme que pour une
mesure gaussienne centrée $ \mu $ sur $ \mathbb{R}^n $ et deux convexes quelconques
$ K, L \subseteq \mathbb{R}^n $ symétriques par rapport à l'origine,
\[ \mu(K \cap L) \;\ge\; \mu(K)\, \mu(L) . \]
Démontrez le cas où $ K $ et $ L $ sont des bandes $ \{ |x_i| \le a_i \} $ portant sur des
ensembles disjoints de coordonnées, vérifiez par calcul direct le cas bidimensionnel
$ K = \{|x_1| \le a\} $, $ L = \{|x_2| \le b\} $ pour un couple corrélé, puis rédigez un compte
rendu précis de l'état des lieux : quels cas particuliers étaient connus (l'inégalité de \v{S}idák de
1967 pour les rectangles), qui a conjecturé l'énoncé général et quand, et comment il a finalement
été démontré.
\end{problem}

\noindent L'inégalité dit une chose sur laquelle un probabiliste parierait sa maison ---
deux contraintes convexes symétriques sont positivement associées sous une gaussienne, de sorte que
satisfaire l'une rend l'autre plus probable --- et elle a résisté six décennies. Un cas particulier
fut proposé par Dunnett et Sobel en 1955 et la conjecture générale énoncée par Das Gupta, Eaton,
Olkin, Perlman, Savage et Sobel en 1972 ; plusieurs démonstrations annoncées se sont effondrées. En
2014, Thomas Royen, statisticien allemand à la retraite travaillant le matin avant le petit
déjeuner, l'a démontrée en quelques pages comme corollaire d'une généralisation aux lois gamma
multivariées, et l'a publiée dans une revue que presque personne ne lisait ; le résultat n'a été
reconnu par la communauté que vers 2017. État en 2026 : démontré, et vérifié indépendamment --- un
rappel que les mathématiques n'ont d'autre gardien qu'un argument correct.

\begin{refs}
\item Contexte : Inégalité de corrélation gaussienne --- Wikipédia (en anglais) (\url{https://en.wikipedia.org/wiki/Gaussian_correlation_inequality})
\item Article : T. Royen, \og{} A simple proof of the Gaussian correlation conjecture extended to multivariate gamma distributions \fg{}, Far East J. Theor. Stat. 48 (2014), 139--145, arXiv:1408.1028 (\url{https://arxiv.org/abs/1408.1028})
\item Article : Z. \v{S}idák, \og{} Rectangular confidence regions for the means of multivariate normal distributions \fg{}, J. Amer. Statist. Assoc. 62 (1967), 626--633
\end{refs}

\bigskip

\begin{problem}[{À quelle vitesse le théorème central limite converge-t-il ? --- Recherche}]
\textbf{PRB-44.} Soient $ X_1, \dots, X_n $ i.i.d. d'espérance $ 0 $, de variance
$ \sigma^2 $ et telles que $ \rho = E|X_1|^3 < \infty $, et soit $ F_n $ la fonction de
répartition de $ (X_1 + \cdots + X_n)/(\sigma\sqrt{n}) $. Le théorème de Berry-Esseen affirme
qu'il existe une constante universelle $ C $ telle que
\[ \sup_{x \in \mathbb{R}} \big| F_n(x) - \Phi(x) \big| \;\le\; \frac{C \rho}{\sigma^3 \sqrt{n}} . \]
\begin{enumerate}
\item[(a)] Démontrez la caractérisation de Stein de la loi normale : une variable aléatoire $ W $ est
normale standard si et seulement si $ E\big[ f'(W) - W f(W) \big] = 0 $ pour toute fonction
$ f $ bornée, continûment dérivable et de dérivée bornée.
\item[(b)] Résolvez l'équation de Stein : étant donnée une fonction test $ h $, montrez que
$ f'(w) - w f(w) = h(w) - E[h(Z)] $ (avec $ Z $ normale standard) admet la solution explicite
\[ f_h(w) = e^{w^2/2} \int_{-\infty}^{w} \big( h(y) - E[h(Z)] \big) e^{-y^2/2} \, dy , \]
et majorez $ f_h $ et $ f_h' $ lorsque $ h $ est l'indicatrice d'une demi-droite.
\item[(c)] Utilisez (a) et (b) pour démontrer une majoration de type Berry-Esseen de la forme annoncée
avec une constante explicite, par un développement de Taylor de
$ E[f'(W_n) - W_n f(W_n)] $ sur la décomposition \og{} une variable en moins \fg{} de la somme.
\item[(d)] Voici maintenant la partie ouverte. Esseen a démontré en 1956 que la meilleure constante
possible vérifie
\[ C \;\ge\; \frac{\sqrt{10}+3}{6\sqrt{2\pi}} \approx 0{,}4097 , \]
au moyen d'une loi à deux points --- reconstruisez un tel exemple et vérifiez la minoration
numériquement. Lisez ensuite l'historique des majorations, depuis le $ 7{,}59 $ d'Esseen lui-même
(1942) en passant par van Beek (1972) et Shiganov (1986) jusqu'aux raffinements de Tiourine et de
Chevtsova vers 2010, qui se situent autour de $ 0{,}47 $ dans le cas i.i.d. et de $ 0{,}56 $
sans loi identique. \textbf{Ouvert :} la constante optimale est inconnue, et la conjecture d'Esseen
selon laquelle sa minoration est la vérité n'est pas démontrée.
\end{enumerate}

\end{problem}

\noindent Andrew Berry (1941) et Carl-Gustav Esseen (1942) ont fourni indépendamment ce
qui manque manifestement au théorème central limite de PRB-14 --- une vitesse --- et leur théorème est
la raison pour laquelle un statisticien peut employer une approximation normale sur un échantillon
de cinquante avec une majoration plutôt qu'un espoir. La voie de démonstration proposée ici est
celle de Charles Stein (1972), inventée parce que les fonctions caractéristiques supportent mal la
dépendance : la méthode de Stein remplace la transformée de Fourier par une caractérisation
différentielle de la loi normale, fonctionne pour des termes dépendants et non identiquement
distribués, et sous-tend aujourd'hui les majorations d'erreur pour l'approximation de Poisson
(Chen), les graphes aléatoires et les algorithmes randomisés. Que quatre-vingts ans d'efforts
n'aient réduit la constante universelle qu'à l'intervalle allant d'environ 0,41 à 0,47 donne une
juste mesure de la délicatesse d'une inégalité optimale. État en 2026 : la constante optimale reste
ouverte.

\begin{refs}
\item Contexte : Inégalité de Berry-Esseen --- Wikipédia (\url{https://fr.wikipedia.org/wiki/In%C3%A9galit%C3%A9_de_Berry-Esseen})
\item Contexte : Méthode de Stein --- Wikipédia (\url{https://fr.wikipedia.org/wiki/M%C3%A9thode_de_Stein})
\item Article : C. Stein, \og{} A bound for the error in the normal approximation to the distribution of a sum of dependent random variables \fg{}, Proc. Sixth Berkeley Symposium on Mathematical Statistics and Probability, vol. 2 (1972)
\item Lecture : Louis H. Y. Chen, Larry Goldstein \& Qi-Man Shao, \textit{Normal Approximation by Stein's Method} (Springer, 2011)
\end{refs}

\bigskip


\vfill
\noindent\emph{Hors de la Caverne} --- des probl\`emes, pas de r\'eponses.
\end{document}
